\documentclass[11pt]{article}
\usepackage[T2A]{fontenc}
\usepackage[utf8]{inputenc}
\usepackage[russian]{babel}
\AtBeginDocument{\shorthandoff{"}}
\usepackage{lmodern}
\usepackage[margin=1in]{geometry}
\usepackage{microtype}

\emergencystretch=2em
\tolerance=1500
\hbadness=10000
\binoppenalty=700
\relpenalty=500

\usepackage{amsmath,amssymb,amsthm,mathtools,bm,mathrsfs}
\usepackage{longtable,booktabs,array,multirow}
\usepackage{textcomp}
\usepackage{graphicx}
\usepackage{enumitem}
\usepackage{xcolor}
\usepackage{framed}
\usepackage{fancyvrb}
\usepackage{ragged2e}
\usepackage{hyperref}
\usepackage{xurl}
\usepackage{tikz}
\usepackage{tikz-cd}
\usetikzlibrary{arrows.meta, positioning, calc, shapes.geometric, fit, decorations.pathreplacing}
\usepackage{caption}

% Theorem environments (Russian)
\theoremstyle{plain}
\newtheorem{theorem}{Теорема}[section]
\newtheorem{lemma}[theorem]{Лемма}
\newtheorem{proposition}[theorem]{Предложение}
\newtheorem{corollary}[theorem]{Следствие}

\theoremstyle{definition}
\newtheorem{definition}[theorem]{Определение}
\newtheorem{example}[theorem]{Пример}
\newtheorem{remark}[theorem]{Замечание}
\newtheorem{convention}[theorem]{Соглашение}

% Custom commands (same as English)
\newcommand{\RS}{\mathrm{R\text{-}S}}
\newcommand{\cS}{\mathcal{S}}
\newcommand{\cF}{\mathcal{F}}
\newcommand{\cM}{\mathcal{M}}
\newcommand{\cC}{\mathcal{C}}
\newcommand{\cD}{\mathcal{D}}
\newcommand{\cL}{\mathcal{L}}
\newcommand{\cP}{\mathcal{P}}
\newcommand{\fM}{\mathfrak{M}}
\newcommand{\Meta}{\mathfrak{Meta}}
\newcommand{\LFnd}{\mathcal{L}_{\mathrm{Fnd}}}
\newcommand{\LCls}{\mathcal{L}_{\mathrm{Cls}}}
\newcommand{\LClsMax}{\mathcal{L}_{\mathrm{Cls}}^{\top}}
\newcommand{\LAbs}{\mathcal{L}_{\mathrm{Abs}}}
\newcommand{\Sint}{\cS_{\mathrm{int}}}
\newcommand{\Ob}{\mathrm{Ob}}
\newcommand{\Mor}{\mathrm{Mor}}
\newcommand{\Hom}{\mathrm{Hom}}
\newcommand{\Fun}{\mathrm{Fun}}
\newcommand{\Cl}{\mathrm{Cl}}
\newcommand{\Lan}{\mathrm{Lan}}
\newcommand{\Ran}{\mathrm{Ran}}
\newcommand{\Set}{\mathrm{Set}}
\newcommand{\Cat}{\mathrm{Cat}}
\newcommand{\Syn}{\mathrm{Syn}}
\newcommand{\Mod}{\mathrm{Mod}}
\newcommand{\Con}{\mathrm{Con}}
\newcommand{\PA}{\mathsf{PA}}
\newcommand{\ZFC}{\mathsf{ZFC}}
\newcommand{\MLTT}{\mathsf{MLTT}}
\newcommand{\HoTT}{\mathsf{HoTT}}
\newcommand{\ETT}{\mathsf{ETT}}
\newcommand{\CIC}{\mathsf{CIC}}
\newcommand{\NBG}{\mathsf{NBG}}
\newcommand{\AFA}{\mathsf{AFA}}
\newcommand{\CZF}{\mathsf{CZF}}
\newcommand{\II}{\mathbf{I}}
\newcommand{\cU}{\mathcal{U}}
\newcommand{\tpi}{\widetilde{\pi}}
\newcommand{\cE}{\mathcal{E}}
\newcommand{\fME}{\fM_{\cE}}
\newcommand{\LAbsE}{\LAbs^{\cE}}
\newcommand{\LClsE}{\LCls^{\cE}}
\newcommand{\LClsMaxE}{(\LClsMax)^{\cE}}
\newcommand{\LFndE}{\LFnd^{\cE}}
\newcommand{\SSE}{\cS_{S}^{\cE}}
\newcommand{\Eff}{\ensuremath{\mathsf{Eff}}}
\newcommand{\NCG}{\ensuremath{\mathsf{NCG}}}
\DeclareMathOperator{\colim}{colim}
\DeclareMathOperator{\image}{image}
\DeclareMathOperator{\id}{id}
\DeclareMathOperator{\ev}{ev}
\DeclareMathOperator{\Trace}{Trace}
\DeclareMathOperator{\dep}{depth}

\newcommand{\gr}{\operatorname{gr}}
\newcommand{\OWL}{\ensuremath{\mathsf{OWL}}}
\newcommand{\SSEbase}{\cS_{S}^{\cE,\,\mathrm{base}}}
\newcommand{\SSEglobal}{\cS_{S}^{\cE,\,\mathrm{global}}}

\title{Пространство модулей формальных систем:\\
Классификация, стабилизация и no-go-теорема для абсолютных оснований}

\author{%
  \footnotesize Макс Середа \\[-3pt]
  \scriptsize\textit{Независимый исследователь} \\[-3pt]
  \scriptsize\texttt{\href{mailto:max@gst.st}{max@gst.st}}%
}

\date{\scriptsize\today}

\begin{document}
\maketitle

\begin{abstract}
Я изучаю $(\infty, n)$-категорное пространство модулей $\fM$ формальных систем (далее \textbf{MSFS}): классифицирующий $2$-стек Морита-классов эквивалентности Rich-оснований ($\LFnd$ — теории, удовлетворяющие условиям (R1)--(R5)). Я строго разграничиваю структурные конструкции, условные критерии сравнения, открытые обязательства стабилизации и относительную граничную лемму.

\emph{(i) Условный плюрализм в $\LCls$.} Классификаторы над общей базой с попарно различными существенными классификационными образами не могут быть эквивалентны над этой базой. Применение критерия к $\infty$-космоям, Унивалентным основаниям и когезивным рамкам требует явного предъявления классификаторов и верификации их образов.

\emph{(ii) Условная мета-категоричность.} Сравнение с общим классификатором даёт биэквивалентность, когда предъявлены существенная сюръективность и эквивалентности на всех hom-категориях. Вывод этих условий поднятия из одних (Max-1)--(Max-4) остаётся открытой задачей.

\emph{(iii) Условное интенсиональное уточнение.} Заданное совместимое сравнение расслаивает уточнение над выбранным базовым образом. Обычная конструкция среза существует, но её display-фильтрация, порождённая одними эквивалентностями, коллапсирует на грейде нуль. Нетривиальное высшее уточнение и строгое Морита-разделение остаются конструктивными обязательствами.

\emph{(iv) Проблема мета-стабилизации.} Эквивалентность между итерированными классификаторами не установлена. Унитарность моделей высших категорий при фиксированном конечном $n$ её не даёт, а неограниченное восхождение по универсумам требует дополнительных допущений.

\emph{(v) Условная AC/OC-двойственность.} До локализации сопряжение $\varepsilon \dashv \alpha$ асимметрично ($\alpha \circ \varepsilon = \id_\cF$, тогда как слой актов $\alpha^{-1}(F)$ удерживает неэквивалентные модели и ресурсные бюджеты); на фиксированном категорном уровне при условии существования локализаций и спуска функторов обращение рефлексионных морфизмов даёт грубую эквивалентность $\fM \simeq \fM_\cE$. Независимо от локализации дуальная граничная лемма $\LAbs^\cE = \emptyset$ выполняется безусловно.

\emph{Граничное следствие: No-go-теорема для абсолютных оснований (AFN-T).} Страта $\LAbs$ (формально определимое, нередуцируемое и максимально порождающее основание) пуста — как по прямому срез-консервативному механизму $(\alpha_0)$, так и по сильной диагональной неабсолютности $(\alpha_{\mathrm{diag}})$ через Фефермана--Пудлака для произвольных подтеорий $\iota_F : F \hookrightarrow S$. AFN-T обобщает классическую no-go-серию Кантор~\cite{Cantor1891}, Рассел~\cite{Russell1903}, Гёдель~\cite{Godel1931}, Тарский~\cite{Tarski1936}, Ловер~\cite{Lawvere1969}, Эрнст~\cite{Ernst2015} и распространяется двойственно на $\LAbs^\cE$.

\smallskip
\noindent\textbf{Ключевые слова:} пространство модулей формальных систем, классифицирующие $2$-стеки, категорная логика, $(\infty, n)$-категории, теория высших топосов, интенсиональная теория типов, мета-классификация, эффективный топос, no-go-теоремы.

\smallskip
\noindent\textbf{MSC2020:} 03A05, 03B30, 03C45, 03F40, 18A05, 18N10, 18N50, 18N60, 18N65.

\end{abstract}

\newpage
\tableofcontents
\newpage

\section{Введение}

\paragraph{Аудит области действия (8 октября 2026).}
Тождественный аргумент AFN-T относителен к определённому классу $\cS_S$ и определению $\Pi_4$, тогда как диагональный аргумент $(\alpha_{\mathrm{diag}})$ исключает внутреннюю интерпретируемость всех Rich-оснований в формально $S$-определимом кандидате. Ни один из них не запрещает новую математику, новые основания или операциональные уточнения Морита-эквивалентных реализаций. В частности, стоимостный или временной инвариант спускается через локализацию лишь тогда, когда он переводит каждый обращаемый морфизм в эквивалентность. Рефлексионная локализация актов не сохраняет автоматически стоимость вычисления, причинность, память или производственные ресурсы; реализация Автогении поэтому требует удержания этих данных \emph{над} локализованным классификатором и отдельных доказательств их сохранения.

Эта статья изучает \emph{пространство модулей формальных систем} $\fM$ — классифицирующий $(\infty, 2)$-стек, точки которого — Морита-классы эквивалентности Rich-оснований (формальных теорий $F$, удовлетворяющих условиям (R1)--(R5) Определения~\ref{def:rs}), а $k$-морфизмы — точные интерпретации и доказуемые естественные эквивалентности. Главные результаты касаются структуры $\fM$:
\begin{itemize}
\item \textbf{Стратификация} (\S\ref{sec:hierarchy}): $\LFnd \xrightarrow{\mathrm{Cls}} \LCls \supseteq \LClsMax \xrightarrow{\mathrm{Gen}} \LAbs = \emptyset$ --- страты высекаются разнотипными пакетами условий ((R1)--(R5) на формальных системах против (M1)--(M5) на классификаторах, Предложение~\ref{prop:level-structure}~(ii)); единственное включение --- $\LClsMax \subseteq \LCls$.
\item \textbf{Условный структурный плюрализм в $\LCls$} (Теорема~\ref{thm:meta-mult}).
\item \textbf{Условная категоричность $\LClsMax$} (Теорема~\ref{thm:meta-cat}).
\item \textbf{Открытые обязательства мета-стабилизации и универсумов} (Замечание~\ref{thm:meta-stab}).
\item \textbf{Условная срез-локальность} (Теорема~\ref{thm:slice-locality}) при предъявленном совместимом сравнении.
\item \textbf{Порождённые фильтрации и их предел} (Определение~\ref{def:display-grade}, Предложение~\ref{prop:grading}): функторы, сохраняющие порождающие, не увеличивают грейд; порождение одними эквивалентностями даёт нулевой грейд всюду (Следствие~\ref{cor:grade-collapse}).
\item \textbf{Пятиосевая абсолютность} (\S\ref{sec:five-axis}).
\item \textbf{Двухэтапное AC/OC-сравнение} (Теорема~\ref{thm:ac-oc-duality}) и безусловная дуальная граничная лемма (Теорема~\ref{thm:dual-afnt}).
\item \textbf{Внешняя граница: $\LAbs = \emptyset$} --- No-go-теорема для абсолютных оснований (Теорема~\ref{thm:afnt-alpha}, Теорема~\ref{thm:afnt-beta}).
\end{itemize}

\paragraph{Машинная верификация.} Сопутствующий корпус Verum должен перепроверяться относительно точных формулировок и допущений настоящей ревизии. Текущий препринт не утверждает машинно-проверенного покрытия каждого утверждения; см.\ верификационный леджер в Приложении~\ref{app:verification-ledger}.

\subsection{Обозначения и соглашения}

\begin{convention}\label{conv:zfc-inacc}
Я работаю всюду в $\ZFC$, дополненной двумя недостижимыми кардиналами $\kappa_1 < \kappa_2$, дающими универсумы Гротендика $\mathbf{U}_1 \subset \mathbf{U}_2$, достаточные для применяемой $2$-категорной техники. Категорный аппарат следует Келли~\cite{Kelly1982} для обогащённых и $2$-категорий, Люри~\cite{LurieHTT} для $(\infty, 1)$-категорий, Люри~\cite{LurieHA} для высшей алгебры и Риль--Верити~\cite{RiehlVerity2022} для $\infty$-космоев. Для теории доступных функторов следую Адамеку--Росицкому~\cite{AdamekRosicky1994}.
\end{convention}

\begin{convention}[Ключевые символы]\label{conv:notation}
\begin{itemize}
\item $\cF$ обозначает $2$-категорию \emph{Rich-оснований} ($\LFnd$): объекты — основополагающие теории в смысле, формализованном в \S\ref{sec:rs} ниже (условия (R1)--(R5)); $1$-морфизмы — точные интерпретации; $2$-морфизмы — доказуемые естественные эквивалентности между интерпретациями.
\item $\rho : \cF \to (\infty, n)\text{-}\Cat$ обозначает \emph{классификационный функтор}, сопоставляющий каждому $F \in \cF$ его категорию моделей (или, для теоретико-типовых $F$, синтаксическую $(\infty, n)$-категорию $\Syn(F)$).
\item \emph{Калибровочное преобразование} $\tau : F_1 \to F_2$ — это $1$-морфизм в $\cF$, сопровождаемый выделенной $2$-эквивалентностью $\rho(F_1) \simeq \rho(F_2)$ на модельно-категорном уровне. Два основания \emph{калибровочно эквивалентны}, если между ними существует обратимое калибровочное преобразование. \emph{Связь с Морита-эквивалентностью.} Калибровочная эквивалентность на уровне $\cF$ соответствует, через фактор $\fM := \cF/\text{калибровка}$, Морита-эквивалентности на уровне стека: $F_1, F_2 \in \cF$ калибровочно эквивалентны тогда и только тогда, когда $[F_1] = [F_2]$ в $\fM$, что по построению есть Морита-класс эквивалентности. Всюду в статье «калибровка» (уровень $1$-морфизмов) и «Морита» (уровень стека) именуют одно и то же лежащее в основе отношение эквивалентности на своих категорных стратах, связанные фактор-функтором $\cF \to \fM$; никакого третьего отношения не вводится.
\item \emph{Морита-редуцируемость на $(\infty, n)$-уровне}: кандидат $X$ в $\LAbs$ и любой $Y \in \cS_S$ рассматриваются как заострённые $(\infty, n)$-категорные структуры через каноническое вложение Йонеды в тотальную $(\infty, n)$-категорию $\mathbf{StrCat}_{S, n}$ $S$-структурированных $(\infty, n)$-категорий (Определение~\ref{def:strcat}, сформулированное сразу ниже). Конкретно: $X$, заданный формулой $\phi_X$, соответствует паре $(\Syn(S), [\phi_X])$, где $[\phi_X]$ — класс изоморфизма объекта, определённого $\phi_X$ в синтаксической категории; $Y \in \Ob(\cM) \subseteq \cS_S^{\mathrm{local}}$ соответствует $(\cM, Y)$. Я пишу $X \to_{M, n} Y$ и говорю, что $X$ \emph{Морита-редуцируем на $(\infty, n)$-уровне} к $Y$, если существует $(\infty, n)$-функтор $\Phi : \Syn(S) \to \cM$ в $\mathbf{StrCat}_{S, n}$, переводящий $[\phi_X]$ в класс изоморфизма $Y$, причём ограничение $\Phi |_{[\phi_X]}$ — $(\infty, n)$-эквивалентность на свой образ. Это расширяет классическую кольцевую Морита-эквивалентность (сравнивающую $\Mod(X)$ и $\Mod(Y)$ как $1$-категории), вбирая полную $(\infty, n)$-категорную структуру; оба понятия совпадают, когда $X, Y$ $1$-категорны и $n = 1$. Я использую сокращение «Морита-эквивалентны» для симметричного понятия и «Морита-редуцируем» для направленной версии.
\item $\fM$ обозначает \emph{классифицирующий $2$-стек} над $\cF$: фактор $\fM := \cF / \text{калибровка}$, рассматриваемый как $(\infty, 2)$-стек Морита-классов эквивалентности Rich-оснований. Формально $\fM$ — выпрямление картезианова расслоения $\int \rho \to \mathrm{R\text{-}S}$, слой которого над $S$ — $2$-группоид $\rho^{-1}(S)$ оснований над $S$ по модулю калибровки. Конструкция следует Люри HTT~\cite{LurieHTT} \S 3.2 (выпрямление--развыпрямление картезиановых расслоений).
\item Для $2$-категории $\mathbf{F}$ обозначение $\mathrm{Aut}_2(\mathbf{F})$ — $2$-группа $2$-автоэквивалентностей $\mathbf{F} \xrightarrow{\simeq_2} \mathbf{F}$ (моноидальная по композиции, с обратимыми $1$-морфизмами и когерентными $2$-клетками).
\item $\Sint$ обозначает $2$-категорию display-$2$-категорий (строгое определение отнесено в Приложение~\ref{app:categorical} ниже).
\item Все $2$-категории локально малы и $\kappa_1$-доступны, если иное не оговорено. (Для конкретных Rich-оснований $S$, чья сила непротиворечивости требует второго недостижимого, параметр доступности $\Syn(S)$ обобщается до $\kappa_S \in \{\kappa_1, \kappa_2\}$ согласно клаузе доступности (R5a) определения Rich-метатеории, формализуемого в \S\ref{sec:rs} ниже; последующие $\kappa_1$-утверждения следует читать равномерно как $\kappa_S$-утверждения, где $S$ индексирован конкретной рассматриваемой $\LFnd$-точкой. Все доказательства статьи устойчивы к этому равномерному чтению, поскольку аргументы доступности используют лишь регулярность кардинала, а не его конкретное значение.)
\item \emph{Усечение и сравнение по глубинам}: $\Syn^{(\infty,n)}(S)$ обозначает представление в явно выбранной модели на глубине $n$, когда такое представление и его отображения сравнения предъявлены. Изменение глубины требует функторов с заявленными универсальными свойствами и когерентных переходных отображений; одна лишь ссылка на левое Кан-расширение не конструирует поднятие отдельно взятой синтаксической категории. Обозначение $n=\infty$ требует отдельно специфицированной модели неограниченной глубины. Ни равенство этой модели неспецифицированному копределу, ни стабилизация под обогащением не предполагаются.
\end{itemize}
\end{convention}

\begin{definition}[Тотальная $(\infty, n)$-категория $\mathbf{StrCat}_{S, n}$]\label{def:strcat}
Для $S \in \RS$ и $n \in \mathbb{N} \cup \{\infty\}$ \emph{тотальная $(\infty, n)$-категория $S$-структурированных $(\infty, n)$-категорий}, обозначаемая $\mathbf{StrCat}_{S, n}$, имеет:
\begin{itemize}
\item \textbf{Объекты}: пары $(\cC, \mathrm{str}_\cC)$, где $\cC$ — $(\infty, n)$-категория, либо существенно $\mathbf{U}_1$-малая, либо $\kappa_1$-доступная (два случая покрывают синтаксические категории и категории моделей соответственно; ср.\ (R5a)), а $\mathrm{str}_\cC$ — $S$-модельная структура на $\cC$ (то есть интерпретация $S$-сигнатуры в $\cC$, удовлетворяющая аксиомам $S$).
\item \textbf{$1$-морфизмы}: $S$-сохраняющие $(\infty, n)$-функторы $\Phi : (\cC_1, \mathrm{str}_1) \to (\cC_2, \mathrm{str}_2)$, то есть функторы, коммутирующие с соответствующими $S$-модельными структурами с точностью до когерентного $(\infty, n)$-изоморфизма.
\item \textbf{Высшие морфизмы}: когерентные $(\infty, n)$-естественные преобразования между $S$-сохраняющими функторами.
\end{itemize}
Объекты $\Syn(S)$ и любой модели $\cM \models S$ я вкладываю в $\mathbf{StrCat}_{S, n}$ конструкцией пойнтинга $X \mapsto (\Syn(S), X)$ и $Y \mapsto (\cM, Y)$ соответственно. Морита-редуцируемость, как указано в Соглашении~\ref{conv:notation}, означает существование $S$-сохраняющего $(\infty, n)$-функтора в $\mathbf{StrCat}_{S, n}$ между такими заострёнными структурами, являющегося эквивалентностью на свой образ.
\end{definition}

\subsection{Реестр допущений}\label{sec:assumption-register}

Реестр разграничивает метатеоретическую базу, структурные определения и дополнительные посылки условных результатов. В частности, критерий сравнения не выводит категоричность из одних (Max-1)--(Max-4), а предлагаемая мета-стабилизация открыта. Последующие использования должны проверять посылки цитируемого утверждения, а не заключать о безусловной доказанности по номеру теоремы.

\begin{description}
\item[\textbf{(A1) Основание.}] $\ZFC$ плюс два сильно недостижимых кардинала $\kappa_1<\kappa_2$ (Соглашение~\ref{conv:zfc-inacc}). \emph{Статус: базовый пакет аксиом.} Это заявленная база для конструкций, условия размера и существования которых верифицированы. Она не удостоверяет собственную непротиворечивость: построение модели даёт лишь верхнюю границу силы непротиворечивости, а равносогласованность требует обратной интерпретации или отдельной нижней оценки.
\item[\textbf{(A2) Скоуп «метатеории».}] No-go квантифицирует по $\RS$, очерченному условиями (R1)--(R5) (\S\ref{sec:rs}): рекурсивно перечислимая, интерпретирующая арифметику Робинсона $Q$, обладающая моделью в $\mathbf{U}_2$, $\mathbf{StrCat}$-совместимая, с существенно $\mathbf{U}_1$-малой $\Syn(S)$ с $\kappa_S$-доступным $\mathrm{Ind}$-пополнением (R5a). \emph{Статус: определенческий скоуп.} Системы вне $\RS$ --- не-р.п.\ оракульные теории, инфинитарные логики без доступного синтаксиса --- не подпадают под вердикт. Теорема~\ref{thm:afnt-alpha} универсальна над $\RS$, а не над неформальной тотальностью «всех оснований».
\item[\textbf{(A3) Формальное содержание «абсолютного основания».}] Членство в $\LAbs$ \emph{определено} как конъюнкция $(F_S) \wedge (\Pi_{4, S, n}) \wedge (\Pi_{3\text{-max}, S, n})$ (Определение~\ref{def:level6}). \emph{Статус: допущение интерпретативной адекватности.} AFN-T закрывает ровно эту тройку; философское чтение «абсолютного основания не существует» верно ровно в той мере, в какой тройка захватывает подразумеваемое понятие. Пятиосевой анализ \S\ref{sec:five-axis} укрепляет это чтение ось за осью, но не устраняет его интерпретативного характера.
\item[\textbf{(A4) Отношение сравнения.}] Объекты сравниваются через Морита-редуцируемость: существование $S$-сохраняющего $(\infty, n)$-функтора в $\mathbf{StrCat}_{S, n}$, являющегося эквивалентностью на свой образ, с объектами, вложенными конструкцией пойнтинга (Определение~\ref{def:strcat}). \emph{Статус: структурная конвенция.} Более тонкая или более грубая топология сравнения перечертила бы границы стратов; \S\ref{sec:intensional-grading} измеряет в точности то, что забывает Морита-фактор.
\item[\textbf{(A5) Исчерпываемость осей.}] Пространством параметров побега объявлены пять осей $(S, n, \mu, \xi, \pi)$ из \S\ref{sec:five-axis}. \emph{Статус: определенческий скоуп, не теорема.} Исчерпываемость выполнена по определению того, что статья считает параметром побега; подлинно новая ось потребовала бы новой теоремы, а не противоречила бы существующей.
\item[\textbf{(A6) Аксиомы мета-классификатора.}] Предлагаемый классифицирующий аппарат удовлетворяет (M1)--(M5), а для максимального класса — (Max-1)--(Max-4). \emph{Статус: аксиомы на предлагаемый классификатор.} Клауза совместного разделения в (Max-4) не даёт ни поднятия морфизмов, ни эквивалентностей на hom-категориях. Теорема~\ref{thm:meta-cat} дополнительно требует функторов сравнения, существенно сюръективных и локально являющихся эквивалентностями. Существование этих данных остаётся открытым.
\item[\textbf{(A7) Локализационная семантика эквивалентности.}] AC/OC-двойственность (\S\ref{sec:ac-oc-duality}) читает грубую «эквивалентность» через $\mathrm{gauge}_\cE$ — $2$-локализацию $\cE$ по классу рефлексий $\mathcal{R}_\cE$ (исчисление дробей Пронк, (BF1)--(BF5)): единица рефлексии обращается \emph{локализацией}, а не в самой $\cE$. \emph{Статус: семантическая конвенция} того же рода, что кольцевая Морита-эквивалентность.
\item[\textbf{(A8) Порождающие display-фильтрации.}] Фильтрация требует явно заданного семени порождающих и допустимых операций. \emph{Статус: данные представления.} Прежнее семя из одних эквивалентностей порождает постоянную фильтрацию грейда нуль. Нетривиальные синтаксические или ресурсно-чувствительные порождающие и инвариантность при их переносе — отдельные обязательства.
\end{description}

\begin{remark}[Интерфейс для предельных расширений]\label{rem:register-interface}
Реестр служит интерфейсом для расширений: каждое расширение должно явно указать, какие допущения оно наследует, а какие заменяет. Расширение может требовать более сильной базы; ни верхняя граница, ни равносогласованность не следуют из одного упоминания базы. AFN-T исключает явно определённую конъюнкцию $\LAbs$ в пределах заявленной области определимости; она не доказывает непротиворечивости иных расширений, категоричности максимального классификатора или потолка для будущих математических конструкций.
\end{remark}

\section{Стратифицированная иерархия пространства модулей}\label{sec:hierarchy}

Я стратифицирую $\fM$ на четыре формальные страты, различаемые двумя структурно различными мета-операциями $\mathrm{Cls}$ (горизонтальная, классифицирующая) и $\mathrm{Gen}$ (вертикальная, порождающая). Страты, определяемые своими категорными условиями, таковы:
\[
\LFnd \;\xrightarrow{\mathrm{Cls}}\; \LCls \;\supsetneq\; \LClsMax \;\xrightarrow{\mathrm{Gen}}\; \LAbs.
\]

\begin{figure}[ht]
\centering
\begin{tikzpicture}[
  level/.style={draw, rounded corners=2pt, minimum width=5cm, minimum height=0.8cm, align=center, font=\small},
  bluefill/.style={fill=blue!6},
  violetfill/.style={fill=violet!10},
  emptybox/.style={draw, dashed, rounded corners=2pt, minimum width=5cm, minimum height=0.8cm, align=center, font=\small, fill=red!5},
  novelty/.style={-{Stealth[length=2mm]}, thick},
  horizmeta/.style={-{Stealth[length=2mm]}, thick, blue!60!black, densely dashed},
  vertmeta/.style={-{Stealth[length=2mm]}, thick, red, dashed},
  node distance=0.45cm
]
\node[level, violetfill] (L5) {$\LFnd$ --- Rich-основания (R1)--(R5)};
\node[level, violetfill, above=of L5] (L5plus) {$\LCls$ --- классификаторы (M1)--(M5)};
\node[level, violetfill, above=of L5plus] (L5max) {$\LClsMax$ --- максимальные (Max-1)--(Max-4)};
\node[emptybox, above=of L5max] (L6) {$\LAbs$ --- \emph{пуста по AFN-T}};
\draw[horizmeta] (L5.north) -- (L5plus.south) node[midway, right, font=\scriptsize, blue!60!black] {$\,\mathrm{Cls}$};
\draw[novelty] (L5plus.north) -- (L5max.south) node[midway, right, font=\scriptsize] {$\,\mathrm{max}$};
\draw[vertmeta] (L5max.north) -- (L6.south) node[midway, right, font=\scriptsize, red] {$\,\mathrm{Gen}$};
\draw[red, very thick] ($(L6.north west)+(0.4,-0.2)$) -- ($(L6.south east)+(-0.4,0.2)$);
\draw[red, very thick] ($(L6.north east)+(-0.4,-0.2)$) -- ($(L6.south west)+(0.4,0.2)$);
\end{tikzpicture}
\caption{Предлагаемая формальная стратификация $\fM$. Классификация и максимальное порождение — разнотипные операции. Условия максимальности определяют $\LClsMax\subseteq\LCls$, но сами по себе не доказывают ни категоричности, ни существования максимальных классификаторов, ни строгости включения. Аргумент AFN-T исключает явно определённую конъюнкцию $\LAbs$. Перенос на акт-центричные модули условен на данных сравнения и локализации Теоремы~\ref{thm:ac-oc-duality}, тогда как $\LAbsE = \emptyset$ выполняется безусловно (Теорема~\ref{thm:dual-afnt}).}
\label{fig:hierarchy}
\end{figure}

\subsection{Формальные страты}

\begin{definition}[Стратифицированная иерархия --- формально]\label{def:hierarchy}
Я определяю формальные \emph{страты} как классы математических объектов:

\begin{tabular}{@{}ll@{}}
\toprule
Класс & Критерий членства \\ \midrule
$\LFnd$ & Формальные системы, удовлетворяющие (R1)--(R5) (уточняются в \S\ref{sec:rs} ниже) \\
$\LCls$ & Классификаторы, удовлетворяющие (M1)--(M5) (уточняются в \S\ref{sec:meta} ниже) \\
$\LClsMax$ & Классификаторы, дополнительно удовлетворяющие (Max-1)--(Max-4) (\S\ref{sec:meta} ниже) \\
$\LAbs$ & Объекты, удовлетворяющие $(F_S) \wedge (\Pi_{4, S, n}) \wedge (\Pi_{3\text{-max}, S, n})$ для некоторого $(S, n)$ (\S\ref{sec:conditions} ниже) \\
\bottomrule
\end{tabular}
\end{definition}

\begin{proposition}[Структурные свойства иерархии]\label{prop:level-structure}
\begin{enumerate}
\item[\textbf{(i)}] \emph{Определимость.} Каждая из четырёх страт — (возможно, собственный) класс, определимый в $\ZFC + 2\text{-inacc}$, дополненной $2$-категорной техникой Соглашения~\ref{conv:zfc-inacc}. Для $\LFnd$ условия — (R1)--(R5) (\S\ref{sec:rs} ниже); для $\LCls$ — (M1)--(M5) (\S\ref{sec:meta} ниже); для $\LClsMax$ они усиливаются до (Max-1)--(Max-4) (\S\ref{sec:meta} ниже); $\LAbs$ определяется тройкой $(F_S) \wedge (\Pi_{4, S, n}) \wedge (\Pi_{3\text{-max}, S, n})$ для некоторого $(S, n)$ (\S\ref{sec:conditions} ниже; структурная самопротиворечивость этой тройки, зафиксированная в \S\ref{sec:afnt-alpha}, даёт Теорему~\ref{thm:afnt-alpha}).
\item[\textbf{(ii)}] \emph{Многостратное членство.} $\LFnd$ и $\LCls$ касаются разнотипных данных: формальной теории и специфицированного классификатора. Построенный объект, несущий обе структуры, может удовлетворять обоим пакетам условий одновременно. Никакое такое членство не выводится из одного названия исследовательской программы.
\item[\textbf{(iii)}] \emph{Включение на максимальном слое.} $\LClsMax\subseteq\LCls$. Строгость включения требует явно построенного классификатора, нарушающего хотя бы одно условие максимальности. Теорема~\ref{thm:meta-cat} даёт условный критерий биэквивалентности при дополнительных данных сравнения, а Теорема~\ref{thm:meta-mult} — условный критерий разделения по образам; ни одна из них сама по себе не строит предлагаемых свидетелей.
\item[\textbf{(iv)}] \emph{Пустота абсолютной страты.} $\LAbs = \emptyset$ по Теореме~\ref{thm:afnt-alpha}. Следовательно, $\LFnd \cap \LAbs = \LCls \cap \LAbs = \LClsMax \cap \LAbs = \emptyset$ тривиально.
\end{enumerate}
\end{proposition}

\begin{proof}
\textbf{(i)} Каждый критерий выразим первопорядково; указанные классы определимы.

\textbf{(ii)} Наборы условий касаются непересекающихся аспектов объекта: аксиомы с исчислением выводов против $2$-категорной структуры с классифицирующим функтором. Объект, несущий обе структуры, удовлетворяет обоим критериям одновременно.

\textbf{(iii)} Включение следует из того, что (Max-1)--(Max-4) усиливают (M1)--(M5). Ни строгость включения, ни эквивалентность максимальных классификаторов не следуют из одной этой импликации.

\textbf{(iv)} Немедленно из Теоремы~\ref{thm:afnt-alpha}.
\end{proof}

\subsection{Примеры}

\paragraph{Страта $\LFnd$.} Основания, с характерными чертами:

\begin{longtable}{@{}p{3.8cm}l p{4cm} p{4.6cm}@{}}
\toprule
Основание & Год & Автор & Характеристика \\ \midrule
Z (Цермело) & 1908 & Цермело & Без схемы подстановки \\
ZF/ZFC & 1922 & Цермело, Френкель & Классическая теория множеств \\
ZFC + недостижимые & 1963+ & Гротендик & С универсумами \\
NBG & 1925--40 & фон Нейман, Бернайс, Гёдель & Классы как объекты \\
MK & 1955 & Морс, Келли & Импредикативные классы \\
KP/KPU & 1965 & Крипке--Платек & Допустимые множества \\
ETCS & 1964 & Ловер & Элементарная теория категории множеств \\
ETCC & 1966 & Ловер & Элементарная теория категорий \\
NF; NFU & 1937; 1969 & Куайн; Йенсен & Стратифицированное свёртывание \\
CZF & 1978 & Акцель & Конструктивная ZF \\
IZF & 1970-е & Фридман & Интуиционистская ZF \\
PA & 1889 & Пеано & Первопорядковая арифметика \\
$\mathsf{Z}_2$ & XX в. & Гильберт--Бернайс & Второпорядковая арифметика \\
MLTT & 1984 & Мартин-Лёф & Зависимые типы \\
CIC & 1988--90 & Кокан--Юэ; Кокан--Полен & Индуктивные конструкции \\
HoTT & 2005+ & Аводей, Воеводский & Унивалентность + типы тождества \\
Кубическая HoTT & 2015+ & Коэн--Кокан и др. & Вычислительная унивалентность \\
Универс-полиморфная HoTT & 2013+ & Воеводский & Полиморфизм универсумов \\
Линейная логика + ! & 1987 & Жирар & Ресурсная чувствительность \\
NBG + AFA & 1988 & Акцель & Необоснованные множества \\
Теория $(\infty, 1)$-топосов & 2009 & Люри & $(\infty, 1)$-топосы \\
Некоммутативная геометрия & 1980-е & Конн & Спектральные тройки \\
Когезивный $(\infty, 1)$-топос & 2013 & Шрайбер & Модальности когезии \\
Мотивная $\mathrm{SH}(k)$ & 1999 & Воеводский, Морель & Мотивная гомотопия \\
Реализуемость (Eff) & 1982 & Хайланд & Эффективный топос \\
Конструктивная математика Маркова & 1950-е--1980-е & А.\,А.~Марков, Нагорный & Конструктивный анализ + принцип Маркова \\
Конструктивный анализ Бишопа & 1967 & Бишоп, Бриджес & Бесвыборный конструктивный вещественный анализ \\
Синтетическая дифф.~геометрия & 1970-е & Кок, Ловер & Инфинитезималии как аксиомы \\
Элементарный высший топос & 2021 & Расех & Элементарный $(\infty, 1)$-топос \\
\bottomrule
\end{longtable}

\paragraph{Страта $\LCls$.} Классификаторы (мета-рамки). Условия разбивают этот предлагаемый класс на $\Meta_{\mathrm{Cls}}^{\top}$ и его дополнение. Перечисленные ниже программы мотивируют кандидатов в классификаторы; сами по себе их названия не задают объектов, удовлетворяющих (M1)--(M5). Категоричность требует гипотез сравнения Теоремы~\ref{thm:meta-cat}.

\begin{longtable}{@{}p{3.3cm}l p{3cm} p{4.4cm}c@{}}
\toprule
Мета-рамка & Год & Автор & Классификационный фокус & В $\Meta_{\mathrm{Cls}}^{\top}$? \\ \midrule
$\infty$-космои & 2022 & Риль--Верити & $(\infty, 1)$-категорные теории & Открыто \\
Унивалентные основания & 2010+ & Воеводский & Расширения HoTT & Открыто \\
Когезивная рамка & 2013 & Шрайбер & Когезивные $\infty$-топосы & Открыто \\
Программа Higher Algebra & 2017+ & Люри & Стабильные $\infty$-кат.\ + операды & Открыто \\
Синтетическая математика & 2000-е+ & Тейлор, Шульман & Аксиоматические синтетические основания & Открыто \\
\bottomrule
\end{longtable}

\paragraph{Страта $\LAbs$.} Пуста, по AFN-T.

\subsection{Две мета-операции}\label{sec:indexing-scheme}

\paragraph{Горизонтальная и вертикальная мета.} На страте $\mathcal{L}$ действуют две категорно различные мета-операции:
\begin{itemize}
\item \textbf{Горизонтальная (классифицирующая) мета.} $\mathrm{Cls}(\mathcal{L})$ — класс $2$-категорий $\mathbf{F}$, \emph{классифицирующих}, но не \emph{порождающих} члены $\mathcal{L}$: локально малых, снабжённых доступным классификационным функтором в модули $\mathcal{L}$, непорождающих и параметризованных метатеорией (ср.\ (M1)--(M5)).
\item \textbf{Вертикальная (порождающая) мета.} $\mathrm{Gen}(\mathcal{L})$ — класс рамок, \emph{порождающих} члены $\mathcal{L}$ как внутренние выводы --- кандидаты в абсолютные основания.
\end{itemize}

\begin{proposition}[Неколлапс горизонтальной меты на $\LFnd$]\label{prop:no-collapse}
$\mathrm{Cls}(\LFnd)$ \emph{не} вкладывается в $\LFnd$ и не совпадает с $\LAbs$. Следовательно, это отдельный формальный страт, который я обозначаю $\LCls$.
\end{proposition}

\begin{proof}[Схема доказательства.]
Член $\mathbf{F} \in \mathrm{Cls}(\LFnd)$ — $2$-категория Rich-оснований с доступным классификационным функтором в $\fM$. Сама она не является Rich-основанием (не удовлетворяет (R1)--(R5) как порождающая теория), поэтому $\mathbf{F} \notin \LFnd$. Она не в $\LAbs$, поскольку (M4) требует непорождаемости, тогда как $\LAbs$ требует максимальной порождающей способности $(\Pi_{3\text{-max}})$.
\end{proof}

\paragraph{Горизонтальная стабилизация.} Эквивалентность итерированных классификаторов — открытая проблема (Замечание~\ref{thm:meta-stab}); AFN-T не доказывает отсутствия дальнейших классификаторных структур.

\paragraph{Вертикальный шаг к запретной страте.} $\LAbs \supseteq \mathrm{Gen}(\LFnd)$: член $\LAbs$ был бы максимально порождающим основанием, из которого выводится всё $\LFnd$. Категорно это свежая операция, а не итерация $\mathrm{Cls}$. По AFN-T, $\LAbs = \emptyset$.

\paragraph{Сводка.} Формальная последовательность читается так:
\[
\underbrace{\LFnd}_{\text{(R1)--(R5)}} \; \xrightarrow{\mathrm{Cls}} \; \underbrace{\LCls}_{\text{(M1)--(M5)}} \; \supseteq \; \underbrace{\LClsMax}_{(\mathrm{Max}\text{-1})\text{--}(\mathrm{Max}\text{-4})} \; \xrightarrow{\mathrm{Gen}} \; \underbrace{\LAbs = \emptyset}_{\text{AFN-T}}
\]

\section{Разумные Rich-метатеории}\label{sec:rs}

\subsection{Условия Rich-метатеории}

\begin{definition}[Разумная Rich-метатеория]\label{def:rs}
Теория $S$, сформулированная в первопорядковом мета-синтаксисе, — \textbf{разумная Rich-метатеория} (или \textbf{R-S}), если она удовлетворяет:
\begin{itemize}
\item \textbf{(R1) Арифметическая полнота}: $S$ интерпретирует арифметику Робинсона $\mathsf{Q}$.
\item \textbf{(R2) Рекурсивная перечислимость}: множество аксиом $S$ рекурсивно перечислимо.
\item \textbf{(R3) Существование модели}: $S$ имеет хотя бы одну модель $\cM$ в объемлющем теоретико-множественном мета-универсуме $\mathbf{U}_2$ (внешний универсум Гротендика Соглашения~\ref{conv:zfc-inacc}, достаточно большой, чтобы вместить модели всех перечисленных R-S-примеров, включая $\ZFC$ + недостижимые кардиналы, чьи наименьшие модели имеют размер $\kappa_1$ и потому лежат в $\mathbf{U}_2 \setminus \mathbf{U}_1$).
\item \textbf{(R4) Гёделево кодирование}: существует рекурсивная инъекция $\ulcorner \cdot \urcorner : \mathrm{Formulas}(S) \to \mathbb{N}$, для которой предикаты «$n$ — гёделев номер доказуемой формулы» и «$n$ — гёделев номер правильно построенной формулы» $S$-определимы.
\item \textbf{(R5) Категорная семантика}: существует $n_S \in \mathbb{N} \cup \{\infty\}$, \emph{внутренняя категорная глубина} $S$, такая что
\begin{itemize}
\item[\textbf{(R5a)}] синтаксическая категория $\Syn(S)$ теории $S$ (объекты: контексты; $1$-морфизмы: доказуемые-из-контекста подстановки; высшие морфизмы: доказательства равенства подстановок) — существенно $\mathbf{U}_1$-малая $(\infty, n_S)$-категория, чья категория моделей $\Mod(S) \simeq \Fun^{\mathrm{lex}}(\Syn(S), -)$ $\kappa_S$-доступна --- эквивалентно, чьё $\mathrm{Ind}_{\kappa_S}$-пополнение $\mathrm{Ind}_{\kappa_S}(\Syn(S))$ есть $\kappa_S$-доступная $(\infty, n_S)$-категория (Адамек--Росицкий~\cite{AdamekRosicky1994}, теорема~2.26; все утверждения о доступности $\Syn(S)$ ниже читаются в этом $\mathrm{Ind}$-пополненном смысле, а утверждения о размере относятся к самой $\Syn(S)$), --- где $\kappa_S \in \{\kappa_1, \kappa_2\}$ определяется теорией $S$: для $S$ силы непротиворечивости $\leq \kappa_1$ (например, $\ZFC$, $\PA$, $\HoTT$, $\MLTT$, $\CIC$) положим $\kappa_S = \kappa_1$; для $S$ силы, требующей второго недостижимого (например, $\ZFC$ + недостижимый), положим $\kappa_S = \kappa_2$. Параметр доступности следует силе непротиворечивости теории; все теоремы ниже формулируются равномерно для произвольного $\kappa_S \in \{\kappa_1, \kappa_2\}$, если конкретный случай не оговорён;
\item[\textbf{(R5b)}] синтаксико-семантическое вычисление $\ev : \Syn(S) \to \cM$ для каждой модели $\cM \models S$ есть $S$-сохраняющий $(\infty, n_S)$-функтор, и пара функторов $(\Syn, \Mod)$ образует сопряжение типа Ламбека--Скотта, единица которого на $\Syn(S)$ является $(\infty, n_S)$-эквивалентностью на \emph{начальную $S$-структурированную $(\infty, n_S)$-категорию}, то есть начальный объект $(\infty, n_S + 1)$-категории $S$-структурированных $(\infty, n_S)$-категорий и $S$-сохраняющих $(\infty, n_S)$-функторов между ними.
\end{itemize}
Синтаксико-семантические данные (R5a)--(R5b) выступают требуемыми посылками на каждую теорию и выбранную категорную модель. Ламбек--Скотт~\cite{LambekScott1986} покрывает обычную категорную логику, а цитируемые гомотопические работы дают конкретную теоретико-типовую семантику. Сравнение моделей конечной глубины~\cite{BarwickSchommerPries2021,BergnerRezk2013} само по себе не строит синтаксико-семантического сопряжения для произвольной формальной теории или произвольной глубины; каждое применение обязано предъявлять требуемую функториальность и универсальное свойство.
\emph{Замечание о внутренней категорной глубине $n_S$.} Параметр $n_S$ означает наибольшее $k$, на котором $\Syn(S)$ несёт структурно нетривиальные $k$-морфизмы до того, как все высшие клетки становятся обратимыми. Для первопорядковой $S$ имеем $n_S = 1$. Для теории типов Мартина-Лёфа и HoTT с унивалентностью синтаксическая категория есть \emph{некая} $(\infty, 1)$-категория в стандартной терминологии (Капулкин--Ламсдейн~\cite{KapulkinLumsdaine2021}); обозначение $n_S = \infty$ в этой статье относится к \emph{неограниченной итерации образования типов тождества}, дающей нетривиальное содержание высших равенств на каждом конечном уровне $k \in \mathbb{N}$, то есть отсутствию такого $k$, на котором вся структура становится $(k-1)$-усечённой. Это чтение «неограниченной глубины» — итеративное понятие, и оно \emph{не} совпадает с $(\infty, \infty)$-категорным понятием Резка / Барвика--Шоммер-Приса~\cite{BarwickSchommerPries2021} (где $k$-морфизмы на каждом $k$ \emph{необратимы}). Оба чтения согласуются в поведении AFN-T (Теорема~\ref{thm:vertical}), и ни один аргумент статьи не зависит от этого различия.

Примеры, верифицирующие (R5a)--(R5b): первопорядковые $S$ при $n_S = 1$ (Ламбек--Скотт~\cite{LambekScott1986}); теория типов Мартина-Лёфа с $n_S = \infty$ (Аводей--Уоррен~\cite{AwodeyWarren2009}); HoTT с $n_S = \infty$ (Капулкин--Ламсдейн~\cite{KapulkinLumsdaine2021}); зависимые теории типов (Джейкобс~\cite{Jacobs1999}).
\end{itemize}
\end{definition}

\subsection{Примеры R-S}

Ниже перечислены разумные Rich-метатеории. Каждый пункт — \emph{формальная теория} $S$ (аксиоматическая система с р.п.\ аксиоматизацией (R2), интерпретирующая Q (R1)); категорные структуры вроде топосов выступают как \emph{модели} этих теорий, а не как сами теории, и цитируются через свои внутренние языки.
\begin{itemize}
\item $\ZFC$; $\ZFC$ + недостижимые кардиналы; $\NBG$; Морс--Келли; $\NBG + \AFA$ (Акцель).
\item Интуиционистские теории типов $\MLTT$ и $\CIC$; $\HoTT$ с унивалентностью; универс-полиморфная $\HoTT$.
\item \emph{Линейная логика с экспоненциалом} $!$, снабжённая арифметическими константами, достаточными для интерпретации Q (то есть интуиционистская линейная арифметика $\mathsf{ILL}_{\mathsf{PA}}$ или ограниченный линейно-арифметический фрагмент лудики Жирара~\cite{Girard2001}); чисто пропозициональная линейная логика (R1) \emph{не} удовлетворяет и исключена.
\item \emph{Внутренняя зависимая теория типов $\infty$-топосов Люри} (аксиоматическое представление $(\infty, 1)$-топоса как модели гомотопической теории типов с унивалентными универсумами и высшими модальностями, Люри HTT~\cite{LurieHTT} \S 5--\S 6, через соответствие внутренних языков Шульмана~\cite{Shulman2019} и Аводея--Уоррена~\cite{AwodeyWarren2009}); сам $(\infty, 1)$-топос — \emph{каноническая модель} этой теории, а не отдельная теория.
\item \emph{Внутренний язык когезивного $\infty$-топоса} (Шрайбер~\cite{Schreiber2013}): $\HoTT$, расширенная сопряжёнными тройками модальных операторов $(\Pi \dashv \flat \dashv \sharp)$, соответствующих аксиомам когезии, оформленная как теоретико-типовое расширение.
\item \emph{Мотивно-расширенные основания}: $\ZFC$ (или $\HoTT$), расширенная аксиомами существования мотивной стабильной $\infty$-категории $\mathrm{SH}(k)$ и её $\mathrm{A}^1$-гомотопической инвариантности как формальное Rich-расширение; сама мотивная $\infty$-категория — снова модельно-теоретический конструкт, не самостоятельная теория.
\item \emph{Внутренний язык топоса реализуемости} $\mathrm{Eff}$: интуиционистская арифметика высших порядков с реализуемостью (Хайланд~\cite{Hyland1982}), канонической моделью которой служит $\mathrm{Eff}$; ниже используется сокращение «$\mathrm{Eff}$-теория» для этого внутреннего языка.
\item \emph{Конструктивная математика Маркова} (Марков, Нагорный~\cite{MarkovNagornyj1988}: $\mathsf{HA}$ + принцип Маркова, с категорной семантикой в $\mathrm{Eff}$).
\item \emph{Конструктивный анализ Бишопа}~\cite{Bishop1967,BridgesVita2006} (бесвыборный конструктивный вещественный анализ с полной категорной семантикой в любом топосе с объектом натуральных чисел).
\item \emph{Предикативный анализ Фефермана}~\cite{Feferman1964} (системы силы $\mathsf{ATR}_0$ с предикативным свёртыванием).
\end{itemize}

\subsection{Не-примеры}

\begin{itemize}
\item \emph{Аффинная логика без экспоненциала $!$}: нарушает (R1) --- индукция Пеано требует контракции на гипотезе индукции.
\item \emph{Противоречивые теории}: нарушают (R3).
\item \emph{Второпорядковая арифметика с полной семантикой}: нарушает (R2) --- множество общезначимых формул не рекурсивно перечислимо.
\item \emph{Интуитивные «теории всего»} в физике без формальной аксиоматизации: нарушают (R1) и (R2).
\item \emph{Строгий философский ультрафинитизм} (Есенин-Вольпин~\cite{EseninVolpin1961,EseninVolpin1970}): нарушает (R1) при строжайшем чтении --- тотальность функции следования отвергается для «неосуществимо больших» аргументов, так что полная индукция Пеано недоступна. Ограниченные варианты (см.\ пограничные случаи ниже) принадлежат другому страту.
\end{itemize}

\begin{proposition}[Нумерация сама по себе не является свидетелем независимости]\label{prop:rs-independence}
Если теория допускает эффективное гёделево кодирование, дополнительный выбор нерекурсивной нумерации не опровергает существования эффективного кодирования.
\end{proposition}
\begin{proof}[Аудит свидетеля]
Исходное эффективное кодирование остаётся экзистенциальным свидетелем. Замена нумерации меняет представление, а не экзистенциальное свойство самой теории. Следовательно, прежняя контрмодель для (R4) не устанавливает независимости условий (R1)--(R5).
\end{proof}
Остальные предлагавшиеся свидетели независимости также требуют реконструкции в фиксированной сигнатуре с проверкой всех прочих условий. В частности, арифметика Робинсона Q не имеет схемы индукции, поэтому отсутствие полной индукции само по себе не доказывает неинтерпретируемость Q; а противоречивое синтаксическое представление не даёт автоматически требуемого синтаксико-семантического сопряжения. Полная независимость (R1)--(R5) здесь не устанавливается.

\begin{itemize}
\item \emph{Ограниченная арифметика и исчисления осуществимости.} Системы $\mathsf{I\Delta_0}$, $\mathsf{S}_2^i$ / $\mathsf{T}_2^i$ Басса~\cite{Buss1986}, $\mathsf{V}_0$, полиномиально-вычислимые арифметики и калькулюсы теоретико-доказательственной осуществимости формализуют ограниченные фрагменты ультрафинитистской интуиции. Они удовлетворяют (R2)--(R5), но реализуют (R1) лишь в ограниченной форме: объемлющая арифметика интерпретирует фрагмент $\mathsf{PA}$ с ограниченными кванторами, а не полную $\mathsf{PA}$. Они допускаются в $\RS$ как \emph{слабые Rich-метатеории}, классифицирующие соответствующий под-стек $\fM$ на ограниченной внутренней категорной глубине $n_S$. Их Морита-отношение к стандартным $\RS$-точкам изучается средствами обратной математики~\cite{Simpson2009}.
\item \emph{Строгий ультрафинитизм (традиция Есенина-Вольпина).} Строжайший философский ультрафинитизм, отвергающий тотальность следования для «неосуществимо больших» аргументов, выпадает из $\RS$ при любом стандартном чтении (R1) (выше, Не-примеры). Ограниченно-арифметические и полиномиальные варианты выше дают формальные реконструкции, допускаемые в $\RS$ как слабые Rich-метатеории.
\item \emph{Паранепротиворечивые теории.} Приложение~\ref{app:paraconsistent} формулирует условный критерий переноса: подходящее классическое ядро и трансляция, сохраняющая все три граничных предиката, должны быть предъявлены явно.
\end{itemize}

\subsection{\texorpdfstring{Класс $\cS_S$}{Класс S\_S}}

\begin{definition}[$S$-определимые категорные структуры]\label{def:SS}
Пусть $S \in \RS$. Определим индуктивно:
\begin{align*}
\cS_S^{\mathrm{local}} &:= \{(\cM, Y) : \cM \models S,\ Y \in \Ob(\cM)\} \;\cup\; \{(\Syn(S), X) : X \in \Ob(\Syn(S))\}, \\
\cS_S^{(0)} &:= \cS_S^{\mathrm{local}}, \\
\cS_S^{(k+1)} &:= \cS_S^{(k)} \,\cup\, \mathcal{O}\bigl(\cS_S^{(k)}\bigr),
\end{align*}
где $\mathcal{O}$ — замыкание под следующим фиксированным списком $(\infty, n)$-категорных операций, каждая из которых берёт $\cS_S^{(k)}$-объекты на входе и производит объекты в $\cS_S^{(k+1)}$:
\begin{enumerate}[nosep,label=\textrm{(O\arabic*)}]
\item Левые и правые Кан-расширения вдоль $\cS_S^{(k)}$-морфизмов;
\item Конструкции Гротендика $\int F$ для картезиановых расслоений $F : I \to \cS_S^{(k)}$ с $I \in \cS_S^{(k)}$;
\item Пределы и копределы малых диаграмм (всюду «малость» диаграммы означает: индексирована $(\infty, n)$-категорией, являющейся $\mathbf{U}_2$-малой) со значениями в $\cS_S^{(k)}$;
\item Категории срезов/запятых/противоположные, построенные из $\cS_S^{(k)}$-данных;
\item Производные функторы (в смысле Квиллена) между $\cS_S^{(k)}$-объектами, снабжёнными модельно-категорной структурой (когда применимо; вакуумно на объектах без модельной структуры);
\item Симметрично-моноидальные тензорные произведения и внутренний hom в замкнутых моноидальных $\cS_S^{(k)}$-объектах.
\end{enumerate}
Положим $\cS_S^{\mathrm{global}} := \bigcup_{k < \omega} \cS_S^{(k)}$ и $\cS_S := \cS_S^{\mathrm{local}} \cup \cS_S^{\mathrm{global}}$. Все операции (O1)--(O6) $\kappa_1$-доступны (Адамек--Росицкий~\cite{AdamekRosicky1994}); $\cS_S^{\mathrm{global}}$ $\kappa_1$-доступен и $\mathbf{U}_2$-мал.
\end{definition}

\subsection{\texorpdfstring{Ключевая лемма: $S$-определимость влечёт членство в $\cS_S$}{Ключевая лемма: S-определимость влечёт членство в S\_S}}

\begin{lemma}[Лемма об $S$-определимости]\label{lem:SS-membership}
Если $X$ — формально $S$-определимый объект, то есть $X$ задан формулой $\phi_X \in S$, выражающей его характеристические свойства, то $X \in \cS_S$.
\end{lemma}

\begin{proof}
По Соглашению~\ref{conv:notation} (клауза Морита-редуцируемости) объект $X$, непосредственно представленный формулой $\phi_X \in S$, рассматривается в $\mathbf{StrCat}_{S, n}$ как заострённая структура $(\Syn(S), [\phi_X])$. Поскольку $[\phi_X] \in \Ob(\Syn(S))$, эта пара лежит в $\cS_S^{\mathrm{local}}$ \emph{по первой клаузе Определения~\ref{def:SS}} --- никакой операции замыкания не требуется. Следовательно, $X \in \cS_S^{\mathrm{local}} \subseteq \cS_S$. (Когда же $\phi_X$ представлена внутри собственной под-теории $\iota_F : F \hookrightarrow S$, как в Определении~\ref{def:F}, переведённая формула $\iota_F(\phi_X) \in S$ задаёт $(\Syn(S), [\iota_F(\phi_X)]_S) \in \cS_S^{\mathrm{local}}$, тогда как отождествление $\Syn(F)[\phi_X]_F$ с $(\Syn(S), [\iota_F(\phi_X)]_S)$ как заострённых $(\infty, n)$-категорий требует срез-консервативности $\Syn(\iota_F)$ — свойства, которое в общем случае нарушается для неконсервативных расширений вроде $\PA \hookrightarrow \PA + \Con(\PA)$. Теорема~\ref{thm:afnt-alpha} ниже закрывает как срез-консервативный случай, так и произвольные точные под-теории $\iota_F : F \hookrightarrow S$.)
\end{proof}

\subsection{Стратификация по размеру}\label{sec:size-stratification}

\paragraph{Классы размера внутри $\cS_S$.}
\begin{itemize}
\item \textbf{Малые} (элементы $\mathbf{U}_1$): множества ранга $< \kappa_1$.
\item \textbf{Большие} (элементы $\mathbf{U}_2 \setminus \mathbf{U}_1$): множества ранга в $[\kappa_1, \kappa_2)$.
\item \textbf{За пределами $\mathbf{U}_2$}: множества ранга $\geq \kappa_2$, включая $\mathbf{U}_2$-собственные классы.
\end{itemize}

\paragraph{Кардинальные границы объемлющих структур.}
\begin{itemize}
\item Класс $S$-определимости $\cS_S = \cS_S^{\mathrm{local}} \cup \cS_S^{\mathrm{global}}$ индексирован $\mathbf{U}_2$: $\cS_S^{\mathrm{local}} \subseteq \mathbf{U}_1$ (объекты $\mathbf{U}_1$-размерных моделей) и $\cS_S^{\mathrm{global}} \subseteq \mathbf{U}_2$ (производные конструкции над $\cS_S^{\mathrm{local}}$ через $\mathbf{U}_2$-ограниченные функториальные операции).
\item $|\Ob(\cF)/\!\simeq| \leq 2^{\aleph_0} < \kappa_1$; $|\cF| \leq \kappa_2$.
\item $|\Ob(\mathbf{StrCat}_{S, n})/\!\simeq| \leq 2^{\kappa_1} < \kappa_2$.
\item Любая оценка $|\fM / \!\simeq_2|$ требует верифицированной оценки размера для выбранного представления. Оценки для итерированных классификаторов требуют отдельных конструкций и предположений об универсумах; Замечание~\ref{thm:meta-stab} не поставляет башни $(\kappa_k)$.
\end{itemize}

\section{\texorpdfstring{Условия страты $\LAbs$}{Условия страты L\_Abs}}\label{sec:conditions}

\begin{definition}[$(F_S)$]\label{def:F}
$(\infty, n)$-категория $X$ удовлетворяет $(F_S)$ для $S \in \RS$ тогда и только тогда, когда $\exists F \in \cF$ с:
\begin{itemize}
\item[\textbf{(F-1)}] $\iota_F : F \hookrightarrow S$ (включение под-теории);
\item[\textbf{(F-2)}] $\iota_F^* : \Mod(S) \to \Mod(F)$ (индуцированный модельный редукт);
\item[\textbf{(F-3)}] $\phi_X \in F$ с $\Syn(F)[\phi_X] \simeq_{(\infty, n)} X$.
\end{itemize}
\end{definition}

\begin{definition}[$(\Pi_{4, S, n})$]\label{def:pi4}
$(\Pi_{4, S, n})(X)$ выполняется тогда и только тогда, когда $\forall Y \in \cS_S,\ \forall\ (\infty, n)$-функтора $F : X \to Y$, $F$ не является $(\infty, n)$-эквивалентностью на свой образ.
\end{definition}

\begin{definition}[$(\Pi_{3\text{-max}, S, n})$]\label{def:pi3max}
$(\Pi_{3\text{-max}, S, n})(X)$ выполняется тогда и только тогда, когда для каждого Rich-основания $F' \in \cF$ существует сохраняющий структуру $(\infty, n)$-функтор $\Syn(F') \to X$, эквивалентный на свой образ и реализующий $F'$ как внутренний вывод в любом основании $F$, представляющем $X$ (индуцируя тем самым точную относительную интерпретацию $F' \to F$, параллельную $1$-морфизму $(\phi, \Phi)$ в Определении~\ref{def:level6-cE}).
\end{definition}

\begin{definition}[$\LAbs$]\label{def:level6}
$X \in \LAbs$ тогда и только тогда, когда $X \models (F_S) \wedge (\Pi_{4, S, n}) \wedge (\Pi_{3\text{-max}, S, n})$ для некоторого $(S, n) \in \RS \times (\mathbb{N} \cup \{\infty\})$.
\end{definition}

\section{\texorpdfstring{Граничная лемма: пустота $\LAbs$ ($(\alpha)$-часть AFN-T)}{Граничная лемма: пустота L\_Abs ((alpha)-часть AFN-T)}}\label{sec:afnt-alpha}

\begin{theorem}[Граничная лемма: пустота $\LAbs$; $(\alpha)$-часть AFN-T]\label{thm:afnt-alpha}
Для каждой Rich-метатеории $S \in \RS$ и каждого категорного уровня $n \in \mathbb{N} \cup \{\infty\}$ не существует $(\infty, n)$-категорного объекта $X$, удовлетворяющего тройке $\LAbs$
\[
(F_S) \wedge (\Pi_{4, S, n}) \wedge (\Pi_{3\text{-max}, S, n}).
\]
Конкретно, пусть $X$ удовлетворяет $(F_S)$ через $\iota_F : F \hookrightarrow S$ и $\Syn(F)[\phi_X] \simeq_{(\infty, n)} X$:
\begin{enumerate}[nosep,label=(\roman*)]
\item \emph{Прямой / срез-консервативный механизм $(\alpha_0)$}: когда $F = S$ (прямая $S$-определимость) или $\Syn(\iota_F)$ ограничивается на $[\phi_X]$ до $(\infty, n)$-эквивалентности на $(\Syn(S), [\iota_F(\phi_X)]_S)$, несовместна уже под-пара $(F_S) \wedge (\Pi_{4, S, n})$.
\item \emph{Сильная диагональная анти-абсолютность $(\alpha_{\mathrm{diag}})$}: для \emph{произвольного} точного включения под-теории $\iota_F : F \hookrightarrow S$ в Определении~\ref{def:F} (без предположения срез-консервативности $\Syn(\iota_F)$, и даже если $(\Pi_{4, S, n})$ опущено либо проколото исключением само-эквивалентностей) несовместна порождающая пара $(F_S) \wedge (\Pi_{3\text{-max}, S, n})$.
\end{enumerate}
\end{theorem}

\begin{corollary}[Пустота $\LAbs$]\label{cor:level6-empty-alpha}
$\LAbs = \emptyset$.
\end{corollary}

\begin{proof}
Немедленно из Теоремы~\ref{thm:afnt-alpha} и Определения~\ref{def:level6}.
\end{proof}

\begin{proof}[Доказательство Теоремы~\ref{thm:afnt-alpha}]
Пусть $X$ удовлетворяет $(F_S)$ через $F \in \cF$, $\iota_F : F \hookrightarrow S$ и $\phi_X \in F$ с $e : X \xrightarrow{\simeq_{(\infty, n)}} \Syn(F)[\phi_X]$.

\emph{Механизм $(\alpha_0)$ (срез-консервативный случай).} Пусть $F = S$ либо $\Syn(\iota_F)|_{[\phi_X]} : \Syn(F)[\phi_X] \to (\Syn(S), [\iota_F(\phi_X)]_S)$ есть $(\infty, n)$-эквивалентность на свой образ. По Лемме~\ref{lem:SS-membership}, применённой к $\iota_F(\phi_X) \in S$, цель $Y := (\Syn(S), [\iota_F(\phi_X)]_S)$ лежит в $\cS_S^{\mathrm{local}} \subseteq \cS_S$, а композит $\Syn(\iota_F)|_{[\phi_X]} \circ e : X \to Y$ (при $F = S$ просто $\id_X : X \to X$) есть $(\infty, n)$-эквивалентность на свой образ с целью в $\cS_S$, что нарушает $(\Pi_{4, S, n})$.

\emph{Механизм $(\alpha_{\mathrm{diag}})$ (произвольное точное $\iota_F : F \hookrightarrow S$ через диагонализацию).} По (R3) теория $S$ имеет модель $\cM \in \mathbf{U}_2$ в объемлющей метатеории $\ZFC + 2\text{-inacc}$, а значит, синтаксически непротиворечива: в стандартных натуральных числах $\mathbb{N} = \omega \in \mathbf{U}_1$ не существует конечного вывода $\bot$ из $S$. Определим первопорядковую арифметическую теорию $T_S := \PA + \Con(S)$, где $\Con(S)$ — каноническое $\Pi_1^0$-предложение непротиворечивости для р.п.\ представления $S$, доставляемого (R2) и (R4). Поскольку $\mathbb{N} \models \PA$ и $\Con(S)$ истинно в $\mathbb{N}$ (никакая $\omega$-непротиворечивость $S$ не предполагается), стандартная модель $\mathbb{N} \in \mathbf{U}_1 \subset \mathbf{U}_2$ свидетельствует (R3) для $T_S$; условия (R1), (R2), (R4) и первопорядковое (R5) ($n_{T_S} = 1$, $\kappa_{T_S} = \kappa_1$; Ламбек--Скотт~\cite{LambekScott1986}) выполняются немедленно, так что $T_S \in \cF$.

Применение $(\Pi_{3\text{-max}, S, n})(X)$ к $F' := T_S \in \cF$ даёт точную относительную интерпретацию $\tau : T_S \to F$. Композиция с $\iota_F : F \hookrightarrow S$ даёт относительную интерпретацию $J := \iota_F \circ \tau : \PA + \Con(S) \to S$, а тем самым и интерпретацию $\mathsf{Q} + \Con(S)$ в $S$. Прообраз арифметических предложений вдоль $J$ определяет $S_{\mathrm{arith}} := \{\psi \in \mathrm{Sent}(\mathcal{L}_{\mathsf{Q}}) : S \vdash J(\psi)\}$ — непротиворечивое рекурсивно перечислимое односортное расширение $\mathsf{Q} + \Con(S)$. Поскольку любой вывод $0 = 1$ в $S_{\mathrm{arith}}$ примитивно-рекурсивно переводится через $J$ в противоречие в $S$, формула $\Con(S)$ влечёт $\Con(S_{\mathrm{arith}})$ на определимом по Соловею $(\mathsf{I\Delta_0} + \Omega_1)$-срезе внутри $S_{\mathrm{arith}}$, доставляя относительную интерпретацию $\mathsf{Q} + \Con(S_{\mathrm{arith}})$ в непротиворечивой р.п.\ теории $S_{\mathrm{arith}}$. Это противоречит второй теореме Гёделя о неполноте для относительной интерпретируемости в форме Фефермана--Пудлака (Феферман~\cite{Feferman1960}, теор.~5.6; Пудлак~\cite{Pudlak1985}, теор.~2.1). Следовательно, никакой объект $X$ не удовлетворяет $(F_S) \wedge (\Pi_{3\text{-max}, S, n})$ и тем более тройке $\LAbs$.
\end{proof}

\section{\texorpdfstring{Граничная лемма: невозможность побега через предел ($(\beta)$-часть AFN-T)}{Граничная лемма: невозможность побега через предел ((beta)-часть AFN-T)}}\label{sec:afnt-beta}

\begin{theorem}[Граничная лемма: невозможность побега через предел; $(\beta)$-часть AFN-T]\label{thm:afnt-beta}
Пусть $\lambda$ — ординал с $\mathrm{cf}(\lambda) < \kappa_2$, регулярным кардиналом (так что конфинальная подпоследовательность $\{A_{\kappa_i}\}_{i < \mathrm{cf}(\lambda)}$, по которой вычисляется копредел, индексирована в $\mathbf{U}_2$; это рабочая граница для категорной техники, никаких ограничений на сам $\lambda$, кроме существования такой конфинальности), и пусть $\{A_\kappa\}_{\kappa < \lambda}$ — трансфинитная последовательность объектов, удовлетворяющая:
\begin{enumerate}[label=(\roman*)]
\item Каждый $A_\kappa$ формально специфицирован внутри некоторого Rich-основания $F_\kappa$ над $S$.
\item Последовательность монотонна в подходящем $(\infty, n)$-категорном смысле (существуют когерентные функторы перехода $A_\kappa \to A_{\kappa'}$ для $\kappa < \kappa'$).
\item Предел $A_\infty := \colim_{\kappa < \lambda} A_\kappa$ существует в объемлющей $(\infty, n)$-категории.
\end{enumerate}
Тогда $A_\infty \in \cS_S^{\mathrm{global}}$, и, следовательно, $A_\infty$ Морита-редуцируем к элементу $\cS_S$. В частности, $A_\infty$ не удовлетворяет $(\Pi_{4, S, n})$.
\end{theorem}

\begin{proof}
Каждый $A_\kappa \in \cS_S$ по Лемме~\ref{lem:SS-membership}; $S$-определимость функторов перехода делает $D : \lambda \to \cS_S$ (со значениями в заострённых структурах Определения~\ref{def:SS}) $S$-индексированной диаграммой. По Определению~\ref{def:SS}, $\cS_S^{\mathrm{global}}$ замкнут под конструкциями Гротендика и Кан-расширениями вдоль $S$-определимых морфизмов; следовательно, $A_\infty = \colim D \in \cS_S^{\mathrm{global}}$. Доступность: по стандартным свойствам замкнутости доступных категорий относительно фильтрованных копределов (Адамек--Росицкий~\cite{AdamekRosicky1994}, в частности замкнутость $\kappa$-доступных категорий под $\lambda$-фильтрованными копределами для регулярного $\lambda \leq \kappa$ и сохранение доступности при доступной переиндексации) копредел наследует $\kappa_1$-доступность от диаграммы. Морита-редукция к каноническому модельно-категорному копределу противоречит $(\Pi_4)$.
\end{proof}

\subsection{Башни собственных классов не дают побега}

\begin{proposition}[Башни собственных классов и неравномерные башни: дихотомия $\RS$-определимости копредела]\label{prop:proper-class}
Пусть $\{A_\alpha\}_{\alpha \in I}$ — произвольная направленная или индексированная собственным классом башня объектов над Rich-основаниями, и пусть подразумеваемый предельный объект $A_\infty := \colim_{\alpha \in I} A_\alpha$ существует в объемлющем $(\infty, n)$-категорном мета-универсуме. Относительно класса $\RS$ предельный объект $A_\infty$ удовлетворяет ровно одной из двух альтернатив, и ни в одном из случаев $A_\infty$ не принадлежит $\LAbs$:
\begin{itemize}
\item[\textbf{(a)}] \emph{$\RS$-определимый копредел.} Существует хотя бы одна $S' \in \RS$, для которой $A_\infty$ удовлетворяет $(F_{S'})$ (например, когда башня равномерно порождена рекурсивной схемой в некоторой $S_{\mathrm{tow}} \in \RS$, определяющей $A_\infty$, или когда неравномерная диаграмма коллапсирует к конечно специфицируемому копределу). Тогда для каждой $S' \in \RS$, свидетельствующей $(F_{S'})(A_\infty)$, и каждого $n \in \mathbb{N} \cup \{\infty\}$ Теорема~\ref{thm:afnt-alpha}, применённая в $(S', n)$, опровергает конъюнкцию $(F_{S'})(A_\infty) \wedge (\Pi_{4, S', n})(A_\infty) \wedge (\Pi_{3\text{-max}, S', n})(A_\infty)$.
\item[\textbf{(b)}] \emph{$\RS$-неопределимый копредел.} Для каждой $S' \in \RS$ предел $A_\infty$ нарушает $(F_{S'})$. Тогда для каждой пары $(S', n) \in \RS \times (\mathbb{N} \cup \{\infty\})$ нарушен первый конъюнкт Определения~\ref{def:level6} на $A_\infty$.
\end{itemize}
Следовательно, $A_\infty \notin \LAbs$.
\end{proposition}

\begin{proof}
В объемлющей метатеории $\ZFC + 2\text{-inacc}$ закон исключённого третьего разбирает высказывание $\exists S' \in \RS,\ (F_{S'})(A_\infty)$. В случае~(a) предположим от противного, что $A_\infty \in \LAbs$; по Определению~\ref{def:level6} найдётся пара $(S_0, n_0) \in \RS \times (\mathbb{N} \cup \{\infty\})$, для которой $A_\infty \models (F_{S_0}) \wedge (\Pi_{4, S_0, n_0}) \wedge (\Pi_{3\text{-max}, S_0, n_0})$, что невозможно по Теореме~\ref{thm:afnt-alpha}, применённой непосредственно в $(S_0, n_0)$. В случае~(b) для всех $S' \in \RS$ выполнено $\neg(F_{S'})(A_\infty)$, поэтому свидетеля $(S', n)$ для Определения~\ref{def:level6} не существует. В обоих случаях $A_\infty \notin \LAbs$.
\end{proof}

\paragraph{Контрпример к наивному восстановлению диаграммы (терминальный коллапс).}
Принципиально, что дихотомия выше проводится по самому \emph{предельному объекту} $A_\infty$, а не по рекурсивной равномерности входной диаграммы $\alpha \mapsto A_\alpha$. Формальная спецификация $\phi_{A_\infty}$ копредела в общем случае \emph{не} позволяет восстановить рекурсивную спецификацию породившей его диаграммы. Уже для счётной направленной диаграммы $(\mathbb{N}, \le)$: выберем произвольную невычислимую бинарную последовательность $b : \mathbb{N} \to \{0, 1\}$, положим $A_k$ равным заострённому конечному множеству мощности $1 + b_k$ с отмеченной точкой $*_k$, положим $f_{k,k} := \id_{A_k}$, а при $j < k$ определим $f_{j,k} : A_j \to A_k$ как постоянное отображение со значением $*_k$. Эти переходы строго компонуются ($f_{k,m} \circ f_{j,k} = f_{j,m}$ при $j < k < m$), и каждый элемент $A_k$ переходит в $*_{k+1}$ на следующем шаге, так что $\colim_{k < \omega} A_k \cong \{*\}$ — терминальное одноточечное множество. Здесь диаграмма $(A_k)_{k < \omega}$ не является р.п.\ (она кодирует $b$), однако $A_\infty \cong \{*\}$ имеет конечную константную спецификацию и потому попадает в случай~\textbf{(a)} выше — исключаясь из $\LAbs$ по Теореме~\ref{thm:afnt-alpha} без какой-либо некорректной попытки реконструировать $(b_k)$ из $\phi_{A_\infty}$.

\subsection{Конкретные башни}

\paragraph{Категорная башня.} Конечные категорные уровни допускают сравнение внутри выбранных моделей. Переход к модели неограниченной глубины и доказательство инвариантности относительно дальнейшего обогащения составляют отдельные обязательства (Замечание~\ref{thm:bergner-lurie-stab}). Унитарность на конечных уровнях не даёт этой стабилизации.

\paragraph{Топосная башня.} Обычные топосы и программы высших топосов мотивируют иерархию семантических сред. Предлагаемый предел неограниченной глубины требует предъявления диаграммы, когерентных переходных отображений, оценок размера и доказательства универсального свойства. Теорема унитарности Барвика--Шоммер-Приса~\cite{BarwickSchommerPries2021} для конечных $n$ не определяет стабилизированного $(\infty, \infty)$-топоса.

\paragraph{Спектральная башня.} $A_0 = C^*$-алгебра; $A_1 = $ спектральная тройка (Конн); $A_n = (\infty, n)$-версия; каждый уровень — производная конструкция над предыдущим.

\subsection{Пространство траекторий}

Более тонкий вариант башенного подхода вводит выбор на каждом шаге: $A_{\kappa+1}$ не определяется однозначно по $A_\kappa$, а зависит от структурного выбора. Пространство $\mathbb{T}$ всех возможных траекторий — большой категорный объект: дерево, симплициальное множество или группоид путей.

\begin{proposition}\label{prop:trajectory-space}
Пространство траекторных выборов допускает (по меньшей мере) четыре представления в разных категорных разрешениях --- это \emph{не} взаимно эквивалентные объекты, но каждое захватывает траекторные данные на своём уровне усечения, и \emph{каждое} лежит в $\cS_S^{\mathrm{global}}$:
\begin{itemize}
\item как группоид путей $\Pi_1(\fM)$, где $\fM$ — классифицирующий $2$-стек Rich-оснований;
\item как симплициальное множество — нерв $\mathrm{N}(\fM)$;
\item как дерево выборов — дерево уточняющих выборов над $\fM$;
\item как категория отношений $\mathrm{Rel}(\fM)$.
\end{itemize}
В каком бы разрешении ни было взято $\mathbb{T}$, вывод ниже применим.
\end{proposition}

\begin{proof}[Доказательство (схема)]
Каждое из четырёх представлений — стандартная производная конструкция над $(\infty, 1)$-топосом Rich-оснований. См.\ Люри~\cite{LurieHTT} \S 3.2 для отождествления группоидов путей с симплициальными множествами; Пронк~\cite{Pronk1996} для интерпретации через бикатегорию дробей; остальные эквивалентности классические.
\end{proof}

\begin{corollary}
Предел любого траекторного башенного подхода лежит в $\cS_S^{\mathrm{global}}$ и потому Морита-редуцируем внутри классифицирующего стека $\fM$. Побега в $\LAbs$ не происходит.
\end{corollary}

\section{Пятиосевая абсолютность}\label{sec:five-axis}

Теорема~\ref{thm:afnt-alpha} устанавливает AFN-T для фиксированной пары $(S, n)$. Теперь я устанавливаю робастность теоремы по осям вариации. Рисунок~\ref{fig:five-axis} суммирует пять осей и их структурные связи.

\begin{figure}[ht]
\centering
\begin{tikzpicture}[
  axis/.style={-{Stealth[length=2.5mm]}, thick},
  dep/.style={-{Stealth[length=2mm]}, thin, gray, densely dotted},
  label/.style={font=\footnotesize, align=center, inner sep=2pt},
  center/.style={draw, circle, thick, minimum size=1.8cm, align=center, font=\bfseries\small, fill=violet!8}
]
\node[center] (C) at (0,0) {AFN-T};
\node[label, above=1.9cm of C]          (H) {Горизонтальная \\[1pt] \scriptsize\itshape по $S \in \RS$};
\node[label, below=1.9cm of C]          (V) {Вертикальная \\[1pt] \scriptsize\itshape по $n \in \mathbb{N} \cup \{\infty\}$};
\node[label, left=2.1cm of C]           (M) {Мета-вертикальная \\[1pt] \scriptsize\itshape мета-итерации};
\node[label, right=2.1cm of C]          (L) {Латеральная \\[1pt] \scriptsize\itshape альт.\ упорядочения};
\node[label] (K) at (3.3,2.3) {Полнота \\[1pt] \scriptsize\itshape внутри $\cL\cS$};
\draw[axis] (C) -- (H);
\draw[axis] (C) -- (V);
\draw[axis] (C) -- (M);
\draw[axis] (C) -- (L);
\draw[axis] (C) -- (K);
\draw[dep] (K.west) to[bend right=12] node[pos=0.35, above right=2pt and 6pt, font=\scriptsize, text=gray, inner sep=1pt] {поглощает} (H.east);
\draw[dep] (L.south) to[bend left=20] node[midway, right, font=\scriptsize, text=gray, inner sep=1pt] {$\sim$ вертикальная} (V.east);
\draw[dep] (M.south) to[bend right=20] node[midway, left, font=\scriptsize, text=gray, inner sep=1pt] {$\mu$-замыкание} (V.west);
\end{tikzpicture}
\caption{Пять рассматриваемых осей вариации. Препятствие AFN-T действует на каждом этапе, на котором установлены его гипотезы определимости и категорной семантики. Диаграмма сама по себе не доказывает, что эти оси исчерпывают все расширения или что мета-вертикальная итерация является стабилизированным замыканием конечных категорных уровней: это отдельные конструктивные и скоуповые обязательства.}
\label{fig:five-axis}
\end{figure}

\begin{theorem}[Горизонтальная абсолютность]\label{thm:horizontal}
AFN-T (Теорема~\ref{thm:afnt-alpha}) выполняется для каждой $S \in \RS$ равномерно: существует единое доказательство, работающее для всех Rich-метатеорий, удовлетворяющих (R1)--(R5).
\end{theorem}

\begin{proof}
Оба механизма Теоремы~\ref{thm:afnt-alpha} опираются исключительно на пакет условий (R1)--(R5): механизм $(\alpha_0)$ использует Лемму~\ref{lem:SS-membership} и категорный тождественный морфизм, а диагональный механизм $(\alpha_{\mathrm{diag}})$ использует лишь (R1)--(R4) и первопорядковое (R5) для пробной теории $T_S = \PA + \Con(S)$. Никакие дополнительные данные конкретной $S$ не привлекаются, поэтому доказательство равномерно покрывает весь класс $\RS$.
\end{proof}

\begin{theorem}[Относительное вертикальное исключение]\label{thm:vertical}
На каждом категорном уровне, для которого предъявлены синтаксико-семантический мост и членство в производном классе Теоремы~\ref{thm:afnt-alpha}, соответствующий кандидат не удовлетворяет тройке $\LAbs$ (и нарушает $\Pi_4$ в срез-консервативном случае).
\end{theorem}
\begin{proof}
Тождественный морфизм срез-консервативного кандидата — эквивалентность с целью из производного класса, а диагональная редукция $(\alpha_{\mathrm{diag}})$ действует на уровне интерпретаций независимо от выбранной глубины $n$, коль скоро выполнены указанные гипотезы. Сам этот аргумент не строит синтаксиса высшей глубины и не доказывает стабилизацию неограниченной башни.
\end{proof}

\begin{theorem}[Относительное мета-вертикальное исключение]\label{thm:meta-vertical}
Любая мета-итерация, фактически формализованная как $S$-определимый объект производного класса $\cS_S$, нарушает $\Pi_4$ в этой метатеории.
\end{theorem}
\begin{proof}
Тождественный функтор на объекте — эквивалентность на свой образ и принадлежит числу интерпретаций, запрещённых $\Pi_4$. Это тот же аргумент, что и в механизме $(\alpha_0)$ Теоремы~\ref{thm:afnt-alpha}.
\end{proof}
Это относительное исключение не означает, что итерация не добавляет новой структуры, что башня стабилизируется или что все возможные итерации $S$-определимы.

\begin{theorem}[Латеральная абсолютность]\label{thm:lateral}
Все альтернативные категорные упорядочения редуцируются к $(\infty, n)$-иерархии (или обогащают её) и потому удовлетворяют AFN-T.
\end{theorem}

\begin{proof}
Каждое из перечисленных упорядочений допускает канонический перевод в $(\infty, n)$-$\Cat$ для подходящего $n$:
\begin{itemize}
\item Операды $\to$ симметрично-моноидальные $(\infty, 1)$-категории (Люри HA~\cite{LurieHA} \S 5).
\item Двойные категории $\to$ $(\infty, 2)$-категории (Грандис).
\item Глобулярные/кубические/опетопные $\omega$-категории $\simeq (\infty, \infty)$-категории (Ара--Мальциниотис~\cite{AraMaltsiniotis2017}).
\item Стабильные $(\infty, 1)$-категории $\to (\infty, 1)$-категории через спектрификацию.
\item Фьюжн-категории $\to$ конечные $(\infty, 1)$-категории.
\end{itemize}
В каждом случае перевод сохраняет формальную определимость: объекты альтернативного упорядочения соответствуют объектам $(\infty, n)$-иерархии на нужном уровне. Следовательно, AFN-T применяется равномерно по всем упорядочениям.
\end{proof}

\begin{definition}[Lawvere-скоуп]\label{def:law-scope}
Основополагающая теория $F$ принадлежит \textbf{Lawvere-скоупу} $\cL\cS$, если выполняются следующие структурные свойства (метатеория, в которой строятся $\Syn(F)$ и $\Mod(F)$, произвольна, но фиксирована на протяжении всех четырёх клауз):
\begin{itemize}
\item \textbf{(L1)} $F$ имеет синтаксическую категорию $\Syn(F)$ --- малую категорию контекстов и доказуемых подстановок.
\item \textbf{(L2)} $F$ имеет семантическую $2$-категорию $\Mod(F)$ --- $2$-категорию моделей.
\item \textbf{(L3)} Существует каноническое сопряжение $\Syn \dashv \Mod$ (Ловер~\cite{Lawvere1969}).
\item \textbf{(L4)} $\Mod(F)$ реализуема как $(\infty, n)$-категория для некоторого $n$ (традиция Гротендика--Люри).
\end{itemize}
Членство в $\cL\cS$ — структурное свойство $F$ и не предполагает никакой конкретной R-S. В частности, каждая R-S $S$ сама лежит в $\cL\cS$ (взять тавтологические $\Syn = S$, $\Mod = \Mod(S)$ из (R5)), так что понятие свободно от циркулярности для применения в Теореме~\ref{thm:completeness}.
\end{definition}

\begin{theorem}[Ось полноты, условно на Lawvere-скоуп]\label{thm:completeness}
Пусть $F \in \cL\cS$. Тогда каждая структурная вариация $F$ редуцируется к одной из четырёх осей (горизонтальной, вертикальной, мета-вертикальной, латеральной). В частности, AFN-T пятиосево абсолютна внутри $\cL\cS$.
\end{theorem}

\begin{proof}
Пусть $F \in \cL\cS$. По Определению~\ref{def:law-scope}, $F$ характеризуется четвёркой:
\begin{itemize}
\item \textbf{(A) Синтаксическая структура $\Syn(F)$}: сигнатура + аксиомы + правила вывода --- \emph{горизонтальная ось} (по R-S).
\item \textbf{(B) Семантическая глубина}: уровень $n$ категории $\Mod(F)$ как $(\infty, n)$-категории --- \emph{вертикальная ось}.
\item \textbf{(C) Мета-рефлексия}: итерации гёделева кодирования и принципов рефлексии --- \emph{мета-вертикальная ось}.
\item \textbf{(D) Альтернативные семантики}: операдные, глобулярные, кубические, опетопные и т.\,д.\ --- \emph{латеральная ось}.
\end{itemize}
По Определению~\ref{def:law-scope} членство в $\cL\cS$ \emph{означает}, что $F$ представлена ровно данными (L1)--(L4), то есть четвёркой (A)--(D); именно здесь живёт условность теоремы --- внутри $\cL\cS$ исчерпываемость (A)--(D) выполняется по определению скоупа, а не как дополнительное утверждение. Значит, любая мнимая пятая ось $\eta$ — подкомпонента одной из (A)--(D). Либо $\eta \subseteq$ (A), и тогда AFN-T для $\eta$ редуцируется к горизонтальной абсолютности (Теорема~\ref{thm:horizontal}); либо $\eta \subseteq$ (B) — редукция к вертикальной (Теорема~\ref{thm:vertical}); либо $\eta \subseteq$ (C) — к мета-вертикальной (Теорема~\ref{thm:meta-vertical}); либо $\eta \subseteq$ (D) — к латеральной (Теорема~\ref{thm:lateral}). Независимой пятой оси внутри $\cL\cS$ нет.
\end{proof}

\section{Три стандартных обходных пути}\label{sec:bypass}

\begin{theorem}[Определимые когерентные семейства]\label{thm:universe}
Пусть $D : I \to \Cat$ — явно специфицированная малая диаграмма, определимая в $S$, и пусть её предел существует в выбранном универсуме и входит в производный класс $\cS_S$. Тогда её предел принадлежит $\cS_S$ и нарушает $\Pi_4$ по тождественной интерпретации.
\end{theorem}
\begin{proof}
Членство дано явной посылкой замыкания; заключительное утверждение есть механизм $(\alpha_0)$ Теоремы~\ref{thm:afnt-alpha}, применённый к полученному определимому объекту.
\end{proof}
Семейства эндофункторов, индексированные универсумами, не образуют такой диаграммы автоматически: ограничение вдоль $\cU_\ell \hookrightarrow \cU_{\ell'}$ не обязано переводить эндофунктор большего универсума в эндофунктор меньшего. Сохранение меньшего универсума и когерентные переходные отображения должны предъявляться явно. Отдельное совместимое семейство — это сечение, а не вся категория всех сечений.

\paragraph{Рефлексивная башня непротиворечивости.} Второй стандартный обход пытается сбежать от AFN-T через итерированные утверждения непротиворечивости. Полагают
\[ S_{\kappa+1} = S_\kappa + \Con(S_\kappa) \]
и надеются, что предел $S_\infty$ захватывает нечто за пределами отправной точки.

\begin{remark}[Итерации непротиворечивости требуют отдельного анализа силы]\label{thm:reflective}
Для заданной рекурсивно аксиоматизированной теории $S$ добавление $\mathrm{Con}(S)$ есть арифметическое расширение непротиворечивости; его сила не равна равномерно «ровно одному недостижимому кардиналу». Трансфинитные итерации требуют явно заданной системы ординальных обозначений, предельного правила и принципа рефлексии; ни одно из них не фиксируется самой записью $S_{\alpha+1} = S_\alpha + \mathrm{Con}(S_\alpha)$. Прежнее утверждение о границе в один недостижимый кардинал и формула сложения $\mathrm{Con}(S)$ с кардиналом отозваны. На каждой корректно определённой Rich-стадии применяется относительный аргумент AFN-T, но сам он не измеряет теоретико-доказательственную силу стадии и не доказывает обрыв башни.
\end{remark}

Интенсиональные представления могут нести данные, забываемые выбранным фактором. Следующий подраздел отделяет верифицированную обычную конструкцию среза от дополнительных условий, необходимых для нетривиального уточнения MSFS. Универсального интенсионального пополнения здесь не строится.

Начну с определения целевой $2$-категории display-$2$-категорий.

\begin{definition}[Display-$2$-класс]\label{def:display-class}
Пусть $\cC$ — $2$-категория с $2$-пулбэками. \textbf{Display-класс} $\cD \subseteq \Mor_1(\cC)$ — класс $1$-морфизмов, удовлетворяющий:
\begin{itemize}
\item \textbf{(D1) Содержит эквивалентности}: каждая $1$-эквивалентность в $\cC$ лежит в $\cD$.
\item \textbf{(D2) Устойчивость к пулбэкам}: для $d : B \to A$ из $\cD$ и любого $f : A' \to A$ $2$-пулбэк $f^*d : f^*B \to A'$ существует и лежит в $\cD$.
\item \textbf{(D3) Замкнутость под композицией}: если $d_1 : C \to B \in \cD$ и $d_2 : B \to A \in \cD$, то $d_2 \circ d_1 \in \cD$.
\item \textbf{(D4) Перенос вдоль обратимых $2$-клеток}: принадлежность к $\cD$ инвариантна относительно обратимых $2$-клеток между параллельными $1$-морфизмами. Выбранные пулбэки снабжены обычными когерентными изоморфизмами сравнения; обратимость переноса вдоль произвольных необратимых $2$-клеток не предполагается.
\end{itemize}
\end{definition}

\begin{definition}[Display-$2$-категория]\label{def:display-2-cat}
\textbf{Display-$2$-категория} — тройка $(\cC, \cD, \iota)$, где:
\begin{itemize}
\item $\cC$ — локально малая $2$-категория с $2$-пулбэками.
\item $\cD$ — display-класс в $\cC$.
\item Рассматривая $\cD$ как полную под-$2$-категорию выбранной $2$-категории стрелок на выделенных display-морфизмах, $\iota : \cD \hookrightarrow \operatorname{Arr}(\cC)$ есть каноническое вложение (конвенция категории стрелок входит в выбор модели).
\end{itemize}
\end{definition}

\begin{definition}[$2$-категория $\Sint$]\label{def:Sint}
В обычном фрагменте $\Sint$ — $2$-категория категорий с пулбэками и display-классами, функторов, сохраняющих пулбэки и display-классы, и естественных преобразований. Соответствующая высшая целевая категория должна задаваться вместе с уровнем преобразований и данными когерентности; приведённый ниже перечень служит типизированным интерфейсом такого расширения, а не доказательством его существования.
\begin{itemize}
\item \textbf{Объекты}: display-$2$-категории $(\cC, \cD, \iota)$.
\item \textbf{$1$-морфизмы}: функтор $F$, сохраняющий выбранные пулбэки с точностью до когерентного изоморфизма и переводящий display-морфизмы в display-морфизмы; индуцированный функтор стрелок ограничивается до $F_\cD$ с $\operatorname{Arr}(F)\iota_1 \simeq \iota_2 F_\cD$. В высшей версии все функторы и клетки когерентности берутся в фиксированной модели.
\item \textbf{$2$-морфизмы}: естественные $2$-преобразования $\eta : F \Rightarrow G$, ограничения которых $\eta|_\cD : F_\cD \Rightarrow G_\cD$ когерентно согласованы.
\end{itemize}
\end{definition}

\begin{theorem}[Верифицированная обычная конструкция среза]\label{thm:I-existence}
Пусть $\mathcal B$ — локально малая обычная категория с выбранными пулбэками. Для $F \in \mathcal B$ положим $\cC_F = \mathcal B / F$, а за $\cD_F$ примем изоморфизмы в $\cC_F$. Тогда
\[
 \II_{\mathcal B} : \mathcal B^{\mathrm{op}} \longrightarrow \Sint,
 \qquad F \longmapsto (\mathcal B / F, \cD_F)
\]
есть контравариантный псевдофунктор в обычном фрагменте display-категорий. Замена базы вдоль изоморфизма является эквивалентностью.
\end{theorem}
\begin{proof}
Срезы имеют пулбэки. Изоморфизмы устойчивы относительно пулбэков, композиции и изоморфизма стрелок. Пулбэк вдоль $f : F_1 \to F_2$ даёт $f^* : \mathcal B / F_2 \to \mathcal B / F_1$, сохраняющий эти данные. Универсальное свойство пулбэков даёт когерентные естественные изоморфизмы $\id^* \simeq \id$ и $(gf)^* \simeq f^* g^*$, а единственность морфизмов сравнения доставляет их тождества единицы и ассоциативности. Обратный изоморфизм базы задаёт квазиобратный функтор замены базы.
\end{proof}

Этот результат сам по себе не строит нетривиального уточнения на всей $2$-категории $\cF$. Прежние свойства (I-1)--(I-4) удерживаются как явные обязательства расширения:
\begin{enumerate}
\item[\textbf{(I-1)}] Задать высшую целевую категорию и когерентное действие на её $2$-клетках; обычный псевдофунктор среза не устанавливает этого для произвольных необратимых клеток $\cF$.
\item[\textbf{(I-2)}] Задать действие калибровочных преобразований; доказанная конструкция сохраняет эквивалентности самой базы, тогда как выбранная эквивалентность категорий моделей не является автоматически эквивалентностью в базе.
\item[\textbf{(I-3)}] Для доказательства строгого уточнения Морита построить два явно заданных представления с равным Морита-образом и неэквивалентными интенсиональными данными. Прежний пример $\MLTT$/$\ETT$ этого не устанавливает: релевантное сравнение использует интенсиональную теорию типов, расширенную одновременно единственностью доказательств тождества и функциональной экстенсиональностью~\cite{KapulkinLi2025}, а не голую $\MLTT$; кроме того, нормализация, разрешимость конверсии и разрешимость типизации — разные свойства, перенос которых требует сохранения эффективного синтаксиса и правил вывода, чего сам по себе display-класс стрелок не даёт.
\item[\textbf{(I-4)}] Предъявить классификационное представление той же вариантности, что и $\II$, когерентное сравнение после забывания display-данных и доказательство предлагаемой локализационной эквивалентности. Контравариантный $\II$ нельзя компоновать и отождествлять с ковариантным $\rho$ без меняющей вариантность конструкции; лемма Йонеды не влечёт $\cF / F \simeq \rho(F)$, а существование локализации само по себе не доказывает её эквивалентности $\fM$.
\end{enumerate}
Высшая конструкция среза требует соответствующих пулбэков и когерентной переиндексации на выбранной базе; эти гипотезы должны проверяться, а не выводиться из одних слов «Rich-основание».

\begin{theorem}[Условная срез-локальность]\label{thm:slice-locality}
Пусть $\mathcal P$ — фиксированная область и предъявлены функторы или псевдофункторы совместимой вариантности
\[
 I : \mathcal P \to \Sint,\qquad \pi : \mathcal P \to \fM,
 \qquad \tpi : \Sint \to \fM
\]
вместе с когерентной эквивалентностью $\tpi I \simeq \pi$. Тогда каждый $I(F)$ имеет базовый образ, эквивалентный $\pi(F)$. Если соответствующий $2$-пулбэк существует, интенсиональный слой над $m$ имеет вид
\[
 \mathrm{Int}(m) := \{m\} \mathbin{\times}^{h}_{\fM} \Sint.
\]
\end{theorem}
\begin{proof}
Первое утверждение есть вычисление заданного сравнения в точке $F$; второе — определение гомотопического слоя с данными сравнения с $m$. Ни нетривиальность слоя, ни свойство расслоения из одного этого определения не следуют.
\end{proof}
Гомотопический слой в общем случае не сводится к полной подкатегории объектов данного класса изоморфизма: его морфизмы обязаны согласовываться с выбранным сравнением с точкой базы. Существование пулбэков в тотальной категории само по себе не делает $\tpi$ расслоением. Построение сравнения в (I-4), нетривиального слоя в (I-3) и данных картезианова поднятия остаётся отдельной задачей.

\begin{corollary}[Постоянство базового образа при заданном сравнении]\label{cor:slice-level}
В предположениях Теоремы~\ref{thm:slice-locality} забывание интенсиональной компоненты $I(F)$ возвращает $\pi(F)$ с точностью до заданной эквивалентности.
\end{corollary}
\begin{proof}
Это компонента эквивалентности $\tpi I \simeq \pi$ на объекте $F$.
\end{proof}
Постоянство выбранного фактора не означает, что уточнение не добавляет математической, вычислительной или ресурсно-чувствительной структуры.

\paragraph{Конечные наблюдательные факторы.}
Для исполняемого языка тождество представлений сравнивает полные программы вместе с локальными привязками библиотек и селекторами вывода. Конечный наблюдательный фактор дополнительно фиксирует интерпретатор, семантику библиотек, свидетельства допуска и разрешённые продолжения. Равенство полных регистровых трасс с отдельным неопределённым исходом задаёт контекстную конгруэнцию лишь тогда, когда удерживаются все метаданные продолжений и допуска и их обновления уважают это равенство. Исчерпание бюджета имеет статус \texttt{Unknown}, отличный от семантической неопределённости и завершённого опровержения. Эти операциональные различия не строят автоматически высшего интенсионального уточнения и не снимают обязательств (I-1)--(I-4).

\subsection{Интенсиональная градуировка}\label{sec:intensional-grading}

Прежнее порождающее правило начинало с эквивалентностей и применяло операции, сохраняющие эквивалентности, а потому не могло дать нетривиальной меры удалённости от эквивалентности. Содержательная фильтрация требует явно выбранных дополнительных порождающих.

\begin{definition}[Грейд относительно выбранных порождающих]\label{def:display-grade}
В заданной категории или $2$-категории выберем класс-семя $G$ морфизмов, насыщенный относительно допустимых обратимых переносов. Положим $D^{(0)} = G$, а $D^{(n+1)}$ определим как насыщение $D^{(n)}$ относительно обратимых переносов вместе с морфизмами, полученными одним разрешённым пулбэком или одной композицией из морфизмов $D^{(n)}$ (так что обратимые данные сравнения имеют нулевую стоимость грейда). Положим $D = \bigcup_{n < \omega} D^{(n)}$ и
\[
 \gr_G(d) = \min\{n : d \in D^{(n)}\}.
\]
Это грейд отдельного морфизма относительно выбранного семени $G$, а не автоматически грейд всей display-категории или объекта интенсионального слоя.
\end{definition}

\begin{proposition}[Функториальность порождённой фильтрации]\label{prop:grading}
Если функтор сохраняет классы-семена и каждую разрешённую порождающую операцию с точностью до допустимых обратимых переносов, то он сохраняет каждую стадию фильтрации и не увеличивает грейд. Композиция удовлетворяет неравенству
\[
 \gr_G(d_2 d_1) \le \max(\gr_G(d_1), \gr_G(d_2)) + 1.
\]
Если класс-семя состоит из всех эквивалентностей и все порождающие операции сохраняют эквивалентности, то каждая стадия совпадает с семенем и каждый грейд равен нулю.
\end{proposition}
\begin{proof}
Функториальное утверждение следует индукцией по стадии порождения. Оба сомножителя композита доступны на наибольшей из своих стадий, поэтому композит появляется не позднее следующей стадии. Для последнего утверждения сохранение эквивалентностей даёт $D^{(n+1)} \subseteq G$, тогда как возрастающая конструкция всегда содержит $G$.
\end{proof}

\begin{corollary}[Коллапс прежнего минимального грейда]\label{cor:grade-collapse}
Для обычной конструкции среза Теоремы~\ref{thm:I-existence} фильтрация, порождённая из изоморфизмов пулбэками, композицией и изоморфизмом стрелок, постоянна на стадии нуль и не может свидетельствовать строгую интенсиональную иерархию.
\end{corollary}
\begin{proof}
Каждая из трёх операций сохраняет изоморфизмы; применяем Предложение~\ref{prop:grading}.
\end{proof}
Эквивалентность между display-категориями — это морфизм в $\Sint$, тогда как эквивалентность в составе display-данных — стрелка внутри одной подлежащей категории: это объекты разных типов, и класс локализации из (I-4) нельзя отождествлять с классом стрелок нулевого грейда только по общему названию.

\begin{remark}[Оставшаяся интенсиональная конструкция]\label{rem:grading-open}
Нетривиальное уточнение требует порождающих, кодирующих синтаксис, вычисление или ресурсы, подходящего отношения сравнения и доказанного переноса этих данных. Инвариантность относительно эквивалентности, разделение Морита-эквивалентных представлений, строгость последовательных стадий и эквивалентность получаемой локализации со стеком $\fM$ остаются открытыми задачами и не следуют из прежнего аргумента о голой паре $\MLTT$/$\ETT$.
\end{remark}

\section{\texorpdfstring{Мета-классификация $\LCls$}{Мета-классификация классификаторов}}\label{sec:meta}

\subsection{\texorpdfstring{Класс $\Meta_{\mathrm{Cls}}$}{Класс Meta\_Cls}}

\begin{definition}[$\LCls$-мета-структура]\label{def:meta}
Объект $\mathbf{F}$ — \textbf{$\LCls$-мета-структура}, если он удовлетворяет:
\begin{itemize}
\item \textbf{(M1) Локально малая $2$-категория}: $\mathbf{F}$ — локально малая $2$-категория.
\item \textbf{(M2) Классификационный функтор}: есть $2$-функтор $\Cl_{\mathbf{F}} : \mathbf{F} \to \fM$ со значениями в классифицирующем $2$-стеке Rich-оснований; пишем $\fM(\mathbf{F}) := \image(\Cl_{\mathbf{F}}) \subseteq \fM$ для под-стека, высекаемого классификацией $\mathbf{F}$.
\item \textbf{(M3) Доступная рефлексия}: есть доступный эндофунктор $R_{\mathbf{F}} : \mathbf{F} \to \mathbf{F}$.
\item \textbf{(M4) Непорождаемость}: $\mathbf{F}$ не выводит аксиом оснований $F \in \image(\Cl_{\mathbf{F}})$ из своих аксиом; классификация параметрическая, не аксиоматическая.
\item \textbf{(M5) Параметризация метатеорией}: $\mathbf{F}$ определён внутри некоторой $S \in \RS$.
\end{itemize}
Класс таких $\mathbf{F}$ обозначаю $\Meta_{\mathrm{Cls}}$.
\end{definition}

\begin{definition}[Максимальность]\label{def:maximality}
Мета-структура $\mathbf{F} \in \Meta_{\mathrm{Cls}}$ \textbf{максимальна}, если дополнительно:
\begin{itemize}
\item \textbf{(Max-1) Полная классификация}: $\image(\Cl_{\mathbf{F}}) = \fM$.
\item \textbf{(Max-2) Калибровочная полнота}: $\mathrm{Aut}_2(\mathbf{F})$ действует транзитивно на калибровочных классах в $\image(\Cl_{\mathbf{F}})$.
\item \textbf{(Max-3) Стратификация по глубине}: $\mathbf{F}$ допускает \emph{глубинно-индексированную фильтрацию} $\mathbf{F}^{(0)} \hookrightarrow \mathbf{F}^{(1)} \hookrightarrow \cdots \hookrightarrow \mathbf{F}^{(n)} \hookrightarrow \cdots$ доступных $2$-подкатегорий с $\mathbf{F} \simeq_2 \colim_{n < \omega} \mathbf{F}^{(n)}$, снабжённую \emph{функцией глубины} $\dep : \Ob(\mathbf{F}) \to \omega$, сопоставляющей каждому объекту наименьшее $n$, при котором объект лежит в $\mathbf{F}^{(n)}$, и удовлетворяющую следующему условию \emph{предикативности}: для каждого предиката $P$, входящего в определение объекта $A \in \mathbf{F}^{(n)}$ свёртыванием (то есть определение имеет вид «$A :=$ объект, классифицирующий $\{x : P(x)\}$»), всякий объект $Y$, выступающий целью под-выражения в $P$ --- в частности, всякая цель $R_\mathbf{F}$, $\Cl_\mathbf{F}$ или вычислений классификатора подобъектов --- удовлетворяет $\dep(Y) < n$.
\item \textbf{(Max-4) Интенсиональная полнота с разделением}: предлагаемый максимальный классификатор должен поставлять корректно типизированное интенсиональное уточнение, удовлетворяющее требованиям (I-1)--(I-4), сформулированным после Теоремы~\ref{thm:I-existence}, вместе с совместным разделением параллельных $1$-морфизмов его экстенсиональным и интенсиональным сравнениями. Это остаётся аксиомой на искомый объект: само разделение не даёт ни недостающих поднятий морфизмов, ни эквивалентностей на hom-категориях.
\end{itemize}
Класс максимальных мета-структур обозначаю $\Meta_{\mathrm{Cls}}^{\top}$.
\end{definition}

\subsection{Мета-категоричность: недостающее условие поднятия}

\begin{theorem}[Условный критерий сравнения]\label{thm:meta-cat}
Пусть $\mathbf F_1, \mathbf F_2$ — локально малые $2$-категории над одной и той же метатеорией. Предположим, что предъявлены общая $2$-категория $\mathbf K$ и $2$-функторы сравнения $\Phi_i : \mathbf F_i \to \mathbf K$, каждый из которых существенно сюръективен на объектах и индуцирует эквивалентность категорий на каждой hom-категории. Тогда $\mathbf F_1 \simeq_2 \mathbf F_2$.
\end{theorem}
\begin{proof}
Каждый $\Phi_i$ является биэквивалентностью в силу заявленных условий на объектах и hom-категориях; компонуем $\Phi_1$ с псевдообратным к $\Phi_2$. Получаемое сравнение зависит от предъявленных данных, и безусловной каноничности без них не утверждается.
\end{proof}

\begin{remark}[Статус предлагаемого максимального классификатора]
Одни условия (Max-1)--(Max-4) не устанавливают этих данных сравнения. Существенная сюръективность на объектах не поднимает морфизмы: функтор из терминальной категории в $BG$ для нетривиальной группы $G$ существенно сюръективен и точен, но не полон. Транзитивное действие на классах объектов не даёт недостающих поднятий, а разделение $1$-морфизмов не даёт эквивалентностей на hom-категориях. Выпрямление применяется к уже заданной индексированной категории или расслоению и не фабрикует требуемого универсального свойства из произвольного классификационного функтора. Построение общего классификатора и доказательство условий поднятия и когерентности остаются открытыми.
\end{remark}

\subsection{Структурная множественность}

\begin{theorem}[Условная структурная множественность]\label{thm:meta-mult}
Пусть $\mathbf F_1, \mathbf F_2, \mathbf F_3$ — построенные мета-структуры, удовлетворяющие (M1)--(M5), с классификационными функторами $\Cl_i$ в общую базу $\fM$. Если их существенные классификационные образы попарно различны, то они попарно неэквивалентны как классификаторы над $\fM$. Если каждый образ собственный, ни одна из них не удовлетворяет (Max-1).
\end{theorem}
\begin{proof}
Эквивалентность классификаторов над $\fM$ сохраняет существенный образ классификационного функтора, поэтому различные образы исключают такую эквивалентность. Собственный образ нарушает требование существенной сюръективности (Max-1).
\end{proof}
Программы Унивалентных оснований, $\infty$-космоев и когезивных высших топосов мотивируют кандидатов для этого сравнения, однако исследовательская программа сама по себе не является специфицированной локально малой $2$-категорией с классификационным функтором MSFS; проверка (M1)--(M5), построение функторов и доказательство различия существенных образов остаются отдельными обязательствами.

\begin{corollary}[Условный плюрализм в $\LCls$]\label{cor:plurality}
При выполнении конструктивных и образных гипотез Теоремы~\ref{thm:meta-mult} страт $\LCls$ содержит несколько попарно неэквивалентных классификаторов над $\fM$.
\end{corollary}

\subsection{Мета-классификация: конструктивные и стабилизационные обязательства}

\begin{remark}[Открытая проблема мета-классификации]\label{thm:meta-stab}
Запись $\fM^{(\mathrm{Cls})} := \Meta_{\mathrm{Cls}} / \text{мета-калибровка}$ именует предлагаемый фактор, но не доказывает его существования с условиями (M1)--(M5): необходимо построить сайт, данные спуска, оценки размера в универсумах и рефлексию. В частности, классификация мета-структур не является автоматически классификационным функтором со значениями в исходном стеке Rich-оснований — цели имеют разные типы.

Эквивалентность этого фактора со вторым итерированным классификатором требует явных функторов сравнения и обратимых единицы и коединицы. Теорема унитарности Барвика--Шоммер-Приса~\cite{BarwickSchommerPries2021} относится к моделям теории $(\infty, n)$-категорий при фиксированном конечном $n$ и не доказывает идемпотентности операции классификации оснований; эта идемпотентность остаётся открытым обязательством. Неограниченная башня недостижимых универсумов также не вытекает из двух недостижимых кардиналов $\kappa_1 < \kappa_2$.
\end{remark}

\begin{corollary}[Относительное исключение на любой построенной стадии]\label{cor:iteration-closure}
На каждой стадии, фактически построенной в выбранной метатеории и удовлетворяющей гипотезам определимости Теоремы~\ref{thm:afnt-alpha}, определяющий её объект не может одновременно удовлетворять $\Pi_4$ (по тождественной интерпретации на данной стадии, без предположения о стабилизации).
\end{corollary}

\begin{remark}[Замкнутость остаётся недоказанной]\label{cor:meta-closure}
Замкнутость $\Meta_{\mathrm{Cls}}$ относительно мета-классификации не вытекает из предыдущего относительного исключения и требует отдельных конструкций и проверок типов, перечисленных в Замечании~\ref{thm:meta-stab}.
\end{remark}

\section{AC/OC-двойственность и дуальная граничная лемма}\label{sec:ac-oc-duality}

Сравнение теоретико-центричных оснований $\cF$ с заострёнными рефлексивными актами $\cE$ проходит в два структурно различных этапа: (1) \emph{до локализации} пара «секция --- забывание» $\varepsilon \dashv \alpha$ (Определение~\ref{def:enactments}) подлинно асимметрична: $\alpha \circ \varepsilon = \id_\cF$ строго, тогда как каждый слой актов $\alpha^{-1}(F)$ удерживает неизоморфные семантические модели, операциональные реализации и ресурсные бюджеты; (2) \emph{после локализации} по рефлексионным морфизмам $\mathcal{R}_\cE$, при выполнении гипотез существования, спуска и когерентности Теоремы~\ref{thm:ac-oc-duality}, грубые стеки $\fM$ и $\fM_\cE$ становятся эквивалентными на фиксированном категорном уровне. Независимо от этих гипотез локализации, дуальная граничная лемма $\LAbsE = \emptyset$ выполняется безусловно в самой $\cE$ (Теорема~\ref{thm:dual-afnt}).

\begin{figure}[ht]
\centering
\begin{tikzpicture}[
  cat/.style={draw, rounded corners=2pt, minimum width=3.6cm, minimum height=0.85cm, align=center, font=\small},
  stack/.style={draw, rounded corners=2pt, minimum width=3.6cm, minimum height=0.85cm, align=center, font=\small, fill=blue!6},
  strat/.style={draw, rounded corners=1.5pt, minimum width=3.2cm, minimum height=0.55cm, align=center, font=\footnotesize, fill=violet!8},
  emptybox/.style={draw, dashed, rounded corners=1.5pt, minimum width=3.2cm, minimum height=0.55cm, align=center, font=\footnotesize, fill=red!6},
  func/.style={-{Stealth[length=2.2mm]}, thick},
  equiv/.style={thick, double, double distance=1.5pt},
  node distance=0.4cm
]
\node[cat, fill=yellow!12] (F) at (-3.6,3.2) {$\cF$ \\[-1pt] \scriptsize\itshape Rich-основания (R1--R5)};
\node[cat, fill=green!12] (E) at (3.6,3.2) {$\cE$ \\[-1pt] \scriptsize\itshape заострённые рефлексивные акты};
\draw[func] ([yshift=5pt]F.east) -- node[midway, above, font=\footnotesize] {$\varepsilon$} ([yshift=5pt]E.west);
\draw[func] ([yshift=-5pt]E.west) -- node[midway, below, font=\footnotesize] {$\alpha$} ([yshift=-5pt]F.east);
\node[font=\scriptsize] at (0, 3.2) {$\dashv$};

\node[stack] (M) at (-3.6, 1.6) {$\fM$ \\[-1pt] \scriptsize\itshape $= \cF / \mathrm{gauge}$};
\node[stack] (ME) at (3.6, 1.6) {$\fME$ \\[-1pt] \scriptsize\itshape $= \cE / \mathrm{gauge}_\cE$};
\draw[func, dashed, gray!60!black] (F) -- node[midway, left, font=\scriptsize, gray!60!black] {фактор} (M);
\draw[func, dashed, gray!60!black] (E) -- node[midway, right, font=\scriptsize, gray!60!black] {фактор} (ME);
\draw[equiv] (M.east) -- node[midway, above, font=\footnotesize] {$\simeq$ (усл.)} node[midway, below, font=\scriptsize, gray!60!black] {Теор.~\ref{thm:ac-oc-duality}} (ME.west);

\node[strat] (L5F) at (-3.6, 0.3) {$\LFnd$};
\node[strat] (L5E) at (3.6, 0.3) {$\LFndE$};
\node[strat] (LCF) at (-3.6, -0.4) {$\LCls$};
\node[strat] (LCE) at (3.6, -0.4) {$\LClsE$};
\node[strat] (LMF) at (-3.6, -1.1) {$\LClsMax$};
\node[strat] (LME) at (3.6, -1.1) {$\LClsMaxE$};
\node[emptybox] (L6F) at (-3.6, -1.8) {$\LAbs = \emptyset$};
\node[emptybox] (L6E) at (3.6, -1.8) {$\LAbsE = \emptyset$};

\draw[equiv, gray!50!black, shorten >=2pt, shorten <=2pt] (L5F) -- (L5E);
\draw[equiv, gray!50!black, shorten >=2pt, shorten <=2pt] (LCF) -- (LCE);
\draw[equiv, gray!50!black, shorten >=2pt, shorten <=2pt] (LMF) -- (LME);
\draw[equiv, red!60!black, shorten >=2pt, shorten <=2pt] (L6F) -- node[midway, above, font=\scriptsize, red!60!black] {Теор.~\ref{thm:dual-afnt}} (L6E);

\draw[gray!40, densely dotted] (M.south) -- (L5F.north);
\draw[gray!40, densely dotted] (ME.south) -- (L5E.north);
\draw[gray!40, densely dotted] (L5F.south) -- (LCF.north);
\draw[gray!40, densely dotted] (L5E.south) -- (LCE.north);
\draw[gray!40, densely dotted] (LCF.south) -- (LMF.north);
\draw[gray!40, densely dotted] (LCE.south) -- (LME.north);
\draw[gray!40, densely dotted] (LMF.south) -- (L6F.north);
\draw[gray!40, densely dotted] (LME.south) -- (L6E.north);
\end{tikzpicture}
\caption{Двухэтапная архитектура AC/OC. На уровне $2$-категорий до локализации (верхний ряд) асимметричная сопряжённая пара $\varepsilon \dashv \alpha$ (Определение~\ref{def:enactments}) вкладывает теоретико-центричную $\cF$ в акт-центричную $\cE$ как корефлексивную секцию ($\alpha \circ \varepsilon = \id_\cF$), тогда как каждый слой актов $\alpha^{-1}(F)$ несёт неизоморфные реализации и операциональные стоимости, коединица рефлексии которых необратима в самой $\cE$. При переходе к локализованным модулям (средний ряд), когда на заданном категорном уровне выполнены гипотезы локализации и спуска Теоремы~\ref{thm:ac-oc-duality}, обращение $\mathcal{R}_\cE$ стягивает каждый слой $\alpha^{-1}(F)$ на $\varepsilon(F)$ и даёт грубую эквивалентность $\fM \simeq \fME$, при которой инвариантные страты соответствуют друг другу (нижние ряды). Независимо от локализации, $\LAbs = \emptyset$ и $\LAbsE = \emptyset$ выполняются безусловно (Теоремы~\ref{thm:afnt-alpha} и~\ref{thm:dual-afnt}).}
\label{fig:ac-oc-duality}
\end{figure}

\subsection{\texorpdfstring{$2$-категория Rich-актов}{2-категория Rich-актов}}

\begin{definition}[$2$-категория актов $\cE$]\label{def:enactments}
\emph{$2$-категория Rich-актов} $\cE$ имеет:
\begin{itemize}
\item \textbf{Объекты}: четвёрки $(F, \cC, \iota, r)$, называемые \emph{заострёнными рефлексивными актами над $F$}, где $F \in \cF$, $\cC \in \mathbf{StrCat}_{S, n_F}$ (Определение~\ref{def:strcat}) — $\kappa_1$-доступная $S$-модельная категория с $F$-модельной структурой, $\iota : \Syn(F) \to \cC$ — $F$-сохраняющий вполне унивалентный $(\infty, n_F)$-функтор на реплетную под-$(\infty, n_F)$-категорию, и $r : \cC \to \Syn(F)$ — \emph{выбранный} левый сопряжённый (рефлектор) к $\iota$ с единицей $\eta_r : \id_\cC \Rightarrow \iota \circ r$ и коединицей $\varepsilon_r : r \circ \iota \Rightarrow \id_{\Syn(F)}$, удовлетворяющими треугольным тождествам в $\mathbf{StrCat}_{S, n_F}$. Рефлектор $r$ — часть данных; он однозначно определён по $\iota$ с точностью до единственной обратимой $2$-клетки (Риль--Верити~\cite{RiehlVerity2022}, предл.~2.1.11; Адамек--Росицкий~\cite{AdamekRosicky1994}, теорема~1.26).
\item \textbf{$1$-морфизмы} $(\phi, \Phi) : (F_1, \cC_1, \iota_1, r_1) \to (F_2, \cC_2, \iota_2, r_2)$: пара из $1$-морфизма $\phi : F_1 \to F_2$ в $\cF$ (точная интерпретация) и $S$-сохраняющего $(\infty, n)$-функтора $\Phi : \cC_1 \to \cC_2$, снабжённого обратимой $2$-клеткой $\Phi \circ \iota_1 \Rightarrow \iota_2 \circ \Syn(\phi)$;
\item \textbf{$2$-морфизмы}: когерентные естественные преобразования $(\phi, \Phi) \Rrightarrow (\phi', \Phi')$, совместимые со структурными клетками, то есть пары $(\tau, T)$ с $\tau : \phi \Rightarrow \phi'$ в $\cF$ и $T : \Phi \Rrightarrow \Phi'$, удовлетворяющие очевидному пятиугольнику с $\iota_1, \iota_2, \Syn(\tau)$.
\end{itemize}
Имеются канонический забывающий $2$-функтор и каноническая секция:
\[
  \alpha : \cE \to \cF, \quad (F, \cC, \iota, r) \mapsto F; \qquad
  \varepsilon : \cF \to \cE, \quad F \mapsto \bigl(F,\, \Syn(F),\, \id_{\Syn(F)},\, \id_{\Syn(F)}\bigr).
\]
По построению $\alpha \circ \varepsilon = \id_\cF$ строго. Для читаемости я иногда пишу $(F, \cC, \iota)$ вместо $(F, \cC, \iota, r)$, когда рефлектор подразумевается.
\end{definition}

\begin{definition}[Калибровка актов]\label{def:gauge-E}
Пусть $\mathcal{R}_\cE \subseteq \Mor_1(\cE)$ состоит из рефлексионных морфизмов $(\tau, \Phi)$, для которых $\tau$ — калибровочное преобразование в $\cF$, а $r_2 \Phi \iota_1$ — эквивалентность. Если выбранная локализация существует с требуемыми оценками размера в универсумах, положим
\[ \mathrm{gauge}_\cE := \text{класс $2$-локализации, порождённый } \mathcal{R}_\cE, \qquad \fM_\cE := \cE[\mathcal{R}_\cE^{-1}]. \]
Построение через бикатегорию дробей Пронк требует проверки всех её посылок, включая условия на $2$-клетки; одна лишь замкнутость под композицией и пулбэками этой проверки не исчерпывает. Зигзаг обращаемых морфизмов даёт эквивалентные объекты, тогда как обратное утверждение требует дополнительного аргумента в выбранной локализации. Рефлексия $(\id_F, r) : (F, \cC, \iota, r) \to \varepsilon(F)$ — предлагаемый к обращению морфизм; её наличие не делает единицу исходной рефлексии обратимой в самой $\cE$.
\end{definition}

\subsection{Акт-синтаксико-семантическая лемма}

\begin{definition}[Класс $\SSE$]\label{def:SSE}
Для $S \in \RS$ класс $\SSE \subseteq \Ob \cE$ \emph{$S$-определимых заострённых актов} — замыкание базового класса
\[
  \SSEbase = \bigl\{ (F', \cM, \iota_\cM, r_\cM) : \cM \models S,\ F' \in \cF,\ \iota_\cM := \ev_\cM,\ r_\cM := (\ev_\cM)^L \bigr\},
\]
где $\ev_\cM : \Syn(F') \to \cM$ — синтаксико-семантическое вычисление из (R5b), а $(\ev_\cM)^L$ — его левый сопряжённый (Ламбек--Скотт), относительно покомпонентного применения операций Определения~\ref{def:SS}: $F$-компонента — по Определению~\ref{def:SS}; $\cC$-компонента — через забывающий $\mathbf{StrCat}_{S, n} \to (\infty, n)\text{-}\Cat$; $(\iota, r)$-компоненты — по функториальности сопряжений под $S$-сохраняющими $(\infty, n)$-функторами. Под-класс $\SSEglobal \subseteq \SSE$ — соответствующее замыкание $\SSEbase$ только теми операциями Определения~\ref{def:SS}, что определяют $\cS_S^{\mathrm{global}}$ (Кан-расширения, конструкции Гротендика, производные функторы).
\end{definition}

\begin{lemma}[Акт-синтаксико-семантическая лемма]\label{lem:SS-membership-E}
Пусть $(F, \cC, \iota, r) \in \cE$ — заострённый рефлексивный акт, четвёрка которого формально $S$-определима, то есть существуют формулы $\phi_F, \psi_\cC, \chi_\iota, \chi_r$ в некотором $F' \in \cF$ с $F' \hookrightarrow S$, специфицирующие соответственно $F, \cC, \iota, r$, причём треугольные тождества и рефлексивность $\iota$ доказуемы в $F'$. Тогда $(F, \cC, \iota, r) \in \SSEglobal$.
\end{lemma}

\begin{proof}
По условию $F$ формально $S$-определим через $\phi_F$; Лемма~\ref{lem:SS-membership} даёт $F \in \cS_S^{\mathrm{global}}$. Модельная категория $\cC$ специфицирована $\psi_\cC$ внутри $F'$; как $S$-структурированная $(\infty, n)$-категория, подлежащая $(\infty, n)$-категория $\cC$ лежит в $\cS_S^{\mathrm{global}}$ по тому же аргументу, применённому к $\cS_S^{\mathrm{global}}$-замыканию Определения~\ref{def:SS}. Пойнтинг $\iota$ и рефлектор $r$, специфицированные $\chi_\iota$ и $\chi_r$ в $F'$, — морфизмы между $\cS_S^{\mathrm{global}}$-объектами, определимые в $F' \hookrightarrow S$. Сопряжённая пара $(\iota, r)$ реализуется в $\SSEglobal$ следующей явной конструкцией: $\iota$ получается как Кан-расширение (Определение~\ref{def:SS}, O1) вдоль вложения $\Syn(F') \hookrightarrow \cC$, специфицированного $\chi_\iota$; существование \emph{левого} рефлектора $r$ (левого сопряжённого к $\iota$ по Определению~\ref{def:enactments}) следует затем из теоремы существования сопряжённого функтора для доступных функторов между локально $\kappa_S$-представимыми категориями (Адамек--Росицкий~\cite{AdamekRosicky1994}, теорема существования сопряжённого функтора в кополном случае), с переносом треугольных тождеств из $\chi_r$. Собранная покомпонентно, четвёрка $(F, \cC, \iota, r) \in \SSEglobal$.
\end{proof}

\subsection{AC/OC Морита-двойственность}

\begin{theorem}[Условное AC/OC-сравнение после локализации]\label{thm:ac-oc-duality}
Зафиксируем категорный уровень. Предположим, что локализации теорий и актов $\fM$ и $\fM_\cE$ существуют, предлагаемые функторы $\varepsilon, \alpha$ спускаются до функторов $\bar\varepsilon, \bar\alpha$ на локализациях и существуют когерентные естественные эквивалентности
\[
 \bar\alpha\bar\varepsilon \simeq \id_{\fM},\qquad
 \bar\varepsilon\bar\alpha \simeq \id_{\fM_\cE}.
\]
Тогда $\fM$ и $\fM_\cE$ эквивалентны, и все предикаты стратов и мета-операции, инвариантные относительно этих данных сравнения, переносятся вдоль эквивалентности.
\end{theorem}
\begin{proof}
Две предъявленные эквивалентности задают квазиобратные функторы сравнения; страт или мета-операция переносится, когда определяющие её данные коммутируют с этими функторами, что включено в заключительную посылку.
\end{proof}
Для рассматриваемых актов равенство $\alpha\varepsilon = \id$ выполняется строго до локализации, а рефлексионные морфизмы мотивируют второе сравнение. Доказательство существования и размера локализаций, спуска обоих функторов и естественности сравнения на всех морфизмах остаётся отдельным обязательством; пообъектное обращение рефлексионных морфизмов само по себе этих условий не заменяет, а поднятие на неограниченную глубину не следует из сравнений на конечных уровнях (Замечание~\ref{thm:bergner-lurie-stab}).

\begin{corollary}[Условная дефинициальная консервативность]\label{cor:ac-oc-conservativity}
Если аппарат актов и все его структурные отображения введены явными определениями, существование и единственность которых доказаны в фиксированной базовой теории $B$, то присоединение их имён консервативно над $B$ для предложений исходного языка $B$.
\end{corollary}
\begin{proof}
Новые имена элиминируются по определяющим формулам; доказанные в $B$ существование и единственность превращают эту элиминацию в трансляцию выводов в $B$ (одна лишь категорная эквивалентность не даёт синтаксической консервативности).
\end{proof}
Применение следствия требует завершения построения аппарата актов в базе $B$. Само по себе использование метатеории $\ZFC + 2\text{-inacc}$ не доказывает равенства $\Con(\mathrm{MSFS}) = \Con(\ZFC + 2\text{-inacc})$: построение модели даёт верхнюю границу силы непротиворечивости, а совпадающая нижняя граница требует обратной интерпретации.

Эквивалентность Теоремы~\ref{thm:ac-oc-duality} переносит условия $\LAbs$ из \S\ref{sec:conditions} на акт-сторону.

\begin{definition}[Условия $\LAbsE$]\label{def:level6-cE}
Для $S \in \RS$ и $n \in \mathbb{N} \cup \{\infty\}$ заострённый рефлексивный акт $(F, \cC, \iota, r) \in \cE$ \textbf{акт-абсолютен}, если он удовлетворяет тройке:
\begin{itemize}
\item[$(\tilde{F}_S)$] \emph{Акт-определимость}: четвёрка $(F, \cC, \iota, r)$ формально $S$-определима в смысле Леммы~\ref{lem:SS-membership-E}, то есть специфицирована формулами $(\phi_F, \psi_\cC, \chi_\iota, \chi_r)$ в некотором $F' \in \cF$ с $F' \hookrightarrow S$ (это уточняет $(F_S)$, дополнительно специфицируя модельно-категорное представление, пойнтинг и рефлектор).
\item[$(\tilde{\Pi}_{4, S, n})$] \emph{Некоординируемость}: не существует $1$-морфизма $(\phi, \Phi) : (F, \cC, \iota, r) \to (F', \cC', \iota', r')$ в $\cE$ с целевой четвёркой из $\SSE$ (Определение~\ref{def:SSE}), в котором $\Phi$ — $(\infty, n)$-эквивалентность на свой образ.
\item[$(\tilde{\Pi}_{3\text{-max}, S, n})$] \emph{Максимальная восприимчивость}: для каждого акта $(F', \cC', \iota', r') \in \cE$ существует $1$-морфизм $(\phi, \Phi) : (F', \cC', \iota', r') \to (F, \cC, \iota, r)$, реализующий Морита-редукцию на $(\infty, n)$-уровне.
\end{itemize}
\emph{Страта $\LAbsE$} — класс $\mathrm{gauge}_\cE$-классов эквивалентности в $\fM_\cE$, имеющих акт-абсолютного представителя.
\end{definition}

\subsection{Дуальная граничная лемма}

\begin{theorem}[Дуальная граничная лемма; $(\alpha)$-часть Dual-AFN-T]\label{thm:dual-afnt}
Для каждой Rich-метатеории $S \in \RS$ и каждого категорного уровня $n \in \mathbb{N} \cup \{\infty\}$ ни один заострённый рефлексивный акт $(F, \cC, \iota, r) \in \cE$ не удовлетворяет тройке $\LAbsE$ $(\tilde{F}_S) \wedge (\tilde{\Pi}_{4, S, n}) \wedge (\tilde{\Pi}_{3\text{-max}, S, n})$. Уже под-пара $(\tilde{F}_S) \wedge (\tilde{\Pi}_{4, S, n})$ структурно самопротиворечива (в срез-консервативном случае), а порождающая пара $(\tilde{F}_S) \wedge (\tilde{\Pi}_{3\text{-max}, S, n})$ для произвольного точного включения $F' \hookrightarrow S$ опровергается проекцией вдоль $\alpha : \cE \to \cF$ на синтический акт $\varepsilon(T_S)$ и механизмом $(\alpha_{\mathrm{diag}})$ Теоремы~\ref{thm:afnt-alpha}.
\end{theorem}

\begin{proof}
В срез-консервативном случае Лемма~\ref{lem:SS-membership-E} даёт $(F, \cC, \iota, r) \in \SSEglobal \subseteq \SSE$, и тождественный $1$-морфизм $(\id_F, \id_\cC)$ нарушает $(\tilde{\Pi}_{4, S, n})$. Для произвольного точного включения $F' \hookrightarrow S$ применение $(\tilde{\Pi}_{3\text{-max}, S, n})$ к $\varepsilon(T_S) \in \cE$ при $T_S = \PA + \Con(S) \in \cF$ даёт $1$-морфизм $(\phi, \Phi) : \varepsilon(T_S) \to (F, \cC, \iota, r)$, забывающая проекция которого $\alpha(\phi, \Phi) = \phi : T_S \to F$ в композиции с представлением $F$ в $F' \hookrightarrow S$ противоречит теореме Фефермана--Пудлака точно как в доказательстве Теоремы~\ref{thm:afnt-alpha}. Следовательно, $\LAbsE = \emptyset$.
\end{proof}

\begin{corollary}[Пустота $\LAbsE$]\label{cor:level6-empty-cE}
$\LAbsE = \emptyset$.
\end{corollary}

\begin{theorem}[Условный перенос относительных осей]\label{thm:dual-five-axis}
Всякий раз, когда осевое сравнение сохраняет посылки относительной граничной леммы и удовлетворяет условиям спуска и когерентности Теоремы~\ref{thm:ac-oc-duality}, соответствующее относительное исключение переносится на представление актов.
\end{theorem}
\begin{proof}
Переносим предполагаемого свидетеля вдоль заданного сравнения и применяем граничную лемму в целевом представлении; сохранение её предикатов входит в посылки переноса.
\end{proof}
Отсюда само по себе не вытекает общее исключение для произвольных смен универсумов, принципов рефлексии или операциональных уточнений; замкнутая симметрично-моноидальная категория не обязана быть картезианово замкнутой и не поддерживает автоматически декартову диагональ Ловера.

\section{Место в серии no-go}\label{sec:series}

AFN-T занимает позицию в структурной серии фундаментальных no-go-теорем. Все такие результаты инстанциируют одну категорную схему: \emph{система, достаточно выразительная, чтобы сформулировать кандидата в максимальный объект, не может содержать этот объект без противоречия}, — реализуемую самореферентным вложением или диагональной интерпретацией. Каждый специализирует схему своим аспектом максимальности и диагональным шагом:

\begin{center}
\begin{tabular}{@{}l l p{3.2cm} p{5.7cm}@{}}
\toprule
Теорема & Год & Аспект максимальности & Диагональный шаг \\ \midrule
Кантор~\cite{Cantor1891} & 1891 & кардинальный & $x \notin f(x)$ для $f : X \to \cP(X)$ \\
Рассел~\cite{Russell1903} & 1903 & свёртывание & $R = \{x : x \notin x\}$, $R \in R \iff R \notin R$ \\
Гёдель (I, II)~\cite{Godel1931} & 1931 & полнота / непротиворечивость & $G = $ «$G$ недоказуемо» \\
Тарский~\cite{Tarski1936} & 1936 & семантическая истина & «это предложение не истинно» \\
Ловер FP~\cite{Lawvere1969} & 1969 & категорная неподвижная точка & картезианова замкнутость + слабая сюръективность $\Rightarrow$ у каждого $f$ есть неподвижная точка \\
Эрнст~\cite{Ernst2015} & 2015 & неограниченная теория категорий & диагональ Ловера при $n = 1$ под Феферманом (R1)--(R3) \\
AFN-T & --- & фундаментальный (определимость + нередуцируемость + порождаемость) & $(\alpha_0)$: синтаксико-семантический мост (Лемма~\ref{lem:SS-membership}) против $(\Pi_4)$; $(\alpha_{\mathrm{diag}})$: $T_S = \PA + \Con(S)$ через Фефермана--Пудлака против $(\Pi_{3\text{-max}})$ (Теорема~\ref{thm:afnt-alpha}) \\
Dual-AFN-T & --- & акт-абсолютный (определимость + некоординируемость + восприимчивость) & тождественный акт $(\id_F, \id_\cC)$ против $(\tilde{\Pi}_4)$ и проекция $\alpha$ на $\varepsilon(T_S)$ против $(\tilde{\Pi}_{3\text{-max}})$ (Теорема~\ref{thm:dual-afnt}) \\
\bottomrule
\end{tabular}
\end{center}

\subsection{Подчинение}

\begin{theorem}[Подчинение через специализацию (уровень рамки)]\label{thm:subsumption}
Для каждой классической no-go-теоремы $T$ из списка ниже существует тройка $(S_T, n_T, X_T)$ с $S_T \in \RS$, $n_T \in \mathbb{N} \cup \{\infty\}$ и $X_T$ — мнимым максимальным объектом $T$, — такая что (i) $X_T \models (F_{S_T})$; (ii) $X_T$ при гипотезе $T$ удовлетворял бы $(\Pi_{4, S_T, n_T})$; (iii) Теорема~\ref{thm:afnt-alpha}, применённая в $(S_T, n_T)$, даёт структурное противоречие с существованием $X_T$ как $\LAbs$-кандидата. Теорема даёт \emph{рамку}, в которой обитает каждый классический no-go; построчные диагонали (канторовская $|X|<|\cP(X)|$, расселовская $R \in R \leftrightarrow R \notin R$ и т.\,д.) — независимые классические доказательства, уточняющие абстрактный шаблон до конкретного свидетеля $T$, и сами по себе не выводятся из одного механизма $(\alpha_0)$.
\begin{center}\footnotesize
\begin{tabular}{@{}lllcl@{}}
\toprule
$T$ & $S_T$ & $X_T$ & $n_T$ & $(F_{S_T})$-свидетель \\ \midrule
Кантор~\cite{Cantor1891} & $\ZFC$ & универсальный $\cP : \mathbf{Set} \to \mathbf{Set}$ & $1$ & $\phi_{\cP}(X,Y) \equiv Y = \cP(X)$ \\
Рассел~\cite{Russell1903} & $\ZFC$ & множество $\{x : x \notin x\}$ & $1$ & $\phi_R(x) \equiv x \notin x$ \\
Гёдель I~\cite{Godel1931} & $\mathsf{PA}$ & $G$ (недоказуемое предложение) & $1$ & $\phi_G \equiv \neg \mathrm{Prov}(\ulcorner G \urcorner)$ \\
Гёдель II~\cite{Godel1931} & $\mathsf{PA}$ & внутреннее доказательство $\mathrm{Con}(\mathsf{PA})$ & $1$ & $\phi_{\mathrm{Con}}$ через $\mathrm{Prov}$ \\
Тарский~\cite{Tarski1936} & $\mathsf{PA}$ & арифметический предикат истины & $1$ & $T(\phi) \equiv \phi$ \\
Ловер FP~\cite{Lawvere1969} & CCC-расш.\ $\ZFC$ & сл.\ сюръ.\ $X$, $f:X \to X^X$ & $1$ & $\phi_X$ в CCC-синтаксисе \\
Эрнст~\cite{Ernst2015} & Феферман (R1)--(R3) & UCT-объект & $1$ & неогр.\ схема Фефермана \\
\bottomrule
\end{tabular}
\end{center}
В каждой строке AFN-T в $(S_T, n_T)$ даёт $\LAbs(S_T, n_T) = \emptyset$ в аспекте максимальности строки; классическая диагональ $T$ — конкретное свидетельское уточнение, устанавливающее $T$ внутри рамки.
\end{theorem}

\begin{proof}
$X_T$ каждой строки прямо $S_T$-определим перечисленной формулой (i), так что $X_T \in \cS_{S_T}^{\mathrm{local}} \subseteq \cS_{S_T}$ по Лемме~\ref{lem:SS-membership}, и $\id_{X_T} : X_T \to X_T$ нарушает $(\Pi_{4, S_T, n_T})$ по Теореме~\ref{thm:afnt-alpha}, откуда $X_T \notin \LAbs$. Построчные классические диагонали уточняют этот шаблон до конкретного свидетеля каждой строки.
\end{proof}

\section{Следствия и открытые вопросы}\label{sec:consequences}

\subsection{Следствия для программ оснований}

AFN-T налагает равномерную структурную границу на кандидатов в абсолютное основание при пустом $\LAbs$:
\begin{itemize}\setlength{\itemsep}{1pt}
\item \textbf{Унивалентные основания}, высшие топосы и когезивные рамки мотивируют кандидатов в классификаторы; их точное членство в стратах MSFS и классификационные образы требуют построения и проверки условий Теоремы~\ref{thm:meta-mult}.
\item \textbf{Теория высших топосов} поставляет семантические структуры; классифицирующий функтор MSFS есть дополнительные данные, без которых программный вердикт о немаксимальности не выводится автоматически.
\item \textbf{Конденсированная математика} (Клаузен--Шольце~\cite{ClausenScholze2019}), перфектоидная геометрия и $\infty$-космои (Риль--Верити~\cite{RiehlVerity2022}) остаются математическими рамками, чья плодотворность и будущие расширения не ограничены относительным $\Pi_4$-исключением.
\item \textbf{Пруф-ассистенты} (Coq, Lean, Agda) реализуют $\LFnd$-основания ($\CIC$, $\CIC + \mathrm{Univalence}$, $\MLTT$); множественность таких систем отражает множественность $\LFnd$.
\item \textbf{Категорная логика}: сравнения фиксированной глубины организуют заданные представления оснований; сама по себе теорема унитарности при конечном $n$ не налагает верхней границы на дальнейшую классификаторную структуру.
\end{itemize}

\paragraph{Диагностика для заданного кандидата.} Сначала необходимо явно задать формальный объект, его метатеорию и отношение интерпретации. Если кандидат принадлежит производному классу $\cS_S$, его тождественный морфизм нарушает $\Pi_4$ по $(\alpha_0)$; если же кандидат представлен в точной подтеории $\iota_F : F \hookrightarrow S$ и претендует на внутреннюю интерпретацию всех Rich-оснований $(\Pi_{3\text{-max}})$, он опровергается диагональным механизмом $(\alpha_{\mathrm{diag}})$ через $T_S = \PA + \Con(S)$. Отнесение конкретных исследовательских программ к $\LCls \setminus \LClsMax$ дополнительно требует построения классификаторов и проверки существенных образов согласно Теореме~\ref{thm:meta-mult}.

\subsection{Открытые вопросы}

\paragraph{(Q1) Существование и сравнение максимальных мета-рамок.} Построить конкретный классификатор, удовлетворяющий (Max-1)--(Max-4) в явно заданной метатеории, а отдельно — общую целевую категорию сравнения с доказательством существенной сюръективности и эквивалентностей на всех hom-категориях (Теорема~\ref{thm:meta-cat}). Сама по себе непустота не влечёт категоричности, а логическая независимость четырёх клауз максимальности остаётся отдельной задачей.

\paragraph{(Q2) Полнота списка мета-рамок.} Насколько полон список $\{\mathbf{F}_{\mathrm{univalent}},\allowbreak \mathbf{F}_{\mathrm{cosmoi}},\allowbreak \mathbf{F}_{\mathrm{cohesive}},\allowbreak \mathbf{F}_{\mathrm{HA}},\allowbreak \mathbf{F}_{\mathrm{synthetic}}\}$ рассматриваемых кандидатов ($\mathbf{F}_{\mathrm{HA}}$ — программа Higher Algebra Люри~\cite{LurieHA})?

\paragraph{(Q3) Исчерпываемость обходных путей.} Теорема~\ref{thm:universe} и Замечание~\ref{thm:reflective} относятся к явно определимым конструкциям при их посылках замыкания и не исключают произвольных смен универсумов, принципов рефлексии или неклассических оснований; формальная теорема исчерпываемости произвольных расширений остаётся открытой.

\paragraph{(Q4) Минимальность силы непротиворечивости.} Определить минимальную базу для каждой завершённой конструкции отдельно. Поскольку мета-стабилизация в Замечании~\ref{thm:meta-stab} открыта, ни достаточность двух недостижимых кардиналов для неё, ни минимальная сила непротиворечивости не установлены.

\paragraph{(Q5) Под-стек слабых Rich-метатеорий.} Образуют ли ограниченно-арифметические теории ($\mathsf{I\Delta_0}$, $\mathsf{S}_2^i$), удовлетворяющие (R2)--(R5) и ограниченному (R1), отдельный страт $\LFnd^{\mathrm{weak}} \subset \LFnd$, или они Морита-плотны в $\LFnd$ средствами обратной математики~\cite{Simpson2009}?

\appendix

\section{Категорные предварительные сведения}\label{app:categorical}

Здесь я собираю категорные инструменты и стандартные теоремы, привлекаемые в доказательствах выше.

\subsection{\texorpdfstring{$2$-категории}{2-категории}}

\textbf{$2$-категория} $\cC$ состоит из класса объектов $\Ob(\cC)$, малой категории $\Hom_\cC(A, B)$ для каждой пары $A, B \in \Ob(\cC)$ и функторов композиции и тождественных $1$-морфизмов со стандартной $2$-категорной когерентностью. Все $2$-категории в статье локально малы, если не оговорено иное (Келли~\cite{Kelly1982}; Бенабу~\cite{Benabou1967}).

\subsection{\texorpdfstring{$(\infty, n)$-категории}{(infinity,n)-категории}}

В $(\infty, n)$-категории имеются $k$-морфизмы при всех $k$, причём $k$-морфизмы при $k > n$ обратимы. Модели на конечных уровнях (квазикатегории, $n$-категории Сегала, полные пространства Сегала, $\Theta_n$-пространства) сравниваются Ара--Мальциниотисом~\cite{AraMaltsiniotis2017}, Люри~\cite{LurieHTT} и Риль--Верити~\cite{RiehlVerity2022}. Стабилизация неограниченной глубины из этих сравнений конечной глубины сама по себе не следует: запись вида $(\infty, \infty + 1)$ требует выбора модели, а её эквивалентность $(\infty, \infty)$ не получается простым повторением унитарности при фиксированном конечном $n$ (см.\ Замечание~\ref{thm:bergner-lurie-stab}).

\subsection{Доступные функторы}

Функтор $F : \cC \to \cD$ между локально малыми категориями \textbf{$\lambda$-доступен} (для регулярного кардинала $\lambda$), если он сохраняет $\lambda$-фильтрованные копределы; категория \textbf{локально $\lambda$-представима}, если она кополна и имеет малую плотную подкатегорию $\lambda$-представимых объектов (Адамек--Росицкий~\cite{AdamekRosicky1994}).

\begin{lemma}[Представление Лэра--Маккаи--Паре]\label{lem:lair}
Каждая доступная категория $2$-категорно эквивалентна $2$-категории моделей малого скетча; эквивалентно, $2$-категории моделей доступной $2$-теории.
\end{lemma}

\noindent\emph{Ссылки для доказательства.} Лэр~\cite{Lair1981}; Маккаи--Паре~\cite{MakkaiPare1989}, глава~3; Адамек--Росицкий~\cite{AdamekRosicky1994}, глава~2.

\subsection{Категории display-отображений}

\begin{definition}[Категория display-отображений, $1$-категорная версия]\label{def:display-1cat}
\textbf{Категория display-отображений} — пара $(\cC, \cD)$, где $\cC$ — категория с пулбэками, а $\cD \subseteq \Mor(\cC)$ — класс морфизмов, устойчивый под пулбэками и содержащий изоморфизмы (Джейкобс~\cite{Jacobs1999}; $2$-категорное расширение — Гамбино--Гарнер~\cite{GambinoGarner2008}).
\end{definition}

\subsection{Бикатегория дробей}

\begin{lemma}[Бикатегория дробей Пронк]\label{lem:pronk}
Даны $2$-категория $\cC$ и класс $\mathcal{W}$ $1$-морфизмов, удовлетворяющий условиям исчисления дробей Бусфилда--Пронк (BF1)--(BF5) (Пронк~\cite{Pronk1996}); тогда $2$-локализация $\cC[\mathcal{W}^{-1}]$ существует как бикатегория дробей: (BF1) тождества; (BF2) композиция; (BF3) $2$-пулбэк-устойчивость; (BF4) правая $2$-фильтрованность и условия на $2$-клетки; (BF5) насыщенность относительно $2$-эквивалентностей.
\end{lemma}

\subsection{Теорема Ловера о неподвижной точке}

\begin{lemma}[Неподвижная точка Ловера, $2$-категорная]\label{lem:lawvere-2}
В $2$-картезианово замкнутой категории $\cC$ со слабо точечно-сюръективным $1$-морфизмом $e : A \to Y^A$ каждый $f : Y \to Y$ имеет неподвижную точку с точностью до когерентной $2$-эквивалентности (Ловер~\cite{Lawvere1969}; Яновский~\cite{Yanofsky2003}).
\end{lemma}

\begin{lemma}[Условная теорема Ловера в гомотопической категории]\label{lem:lawvere-inf}
Пусть выбранная высше-категорная модель имеет картезианово замкнутую гомотопическую категорию $h\cC$, и пусть $e : A \to Y^A$ слабо точечно-сюръективен в $h\cC$. Тогда каждый эндоморфизм $f : Y \to Y$ имеет неподвижную глобальную точку в $h\cC$.
\end{lemma}
\begin{proof}
Применяем обычный категорный диагональный аргумент в $h\cC$: представляем отображение $a \mapsto f(e(a)(a))$ точкой $a_0$ и вычисляем в $a_0$; точка $y_0 = e(a_0)(a_0)$ удовлетворяет $f(y_0) = y_0$ в $h\cC$.
\end{proof}
Это равенство в гомотопической категории $h\cC$; более сильная когерентная высшая структура неподвижной точки требует отдельного определения и дополнительных данных.

\subsection{Критерий типа Уайтхеда}

\begin{lemma}[Условная реконструкция по башне]\label{lem:whitehead}
Если выбранная модель восстанавливает $\cC_i$ как предел башни $\cC_i^{(n)}$, а функтор $\Phi$ есть предел совместимого семейства эквивалентностей $\Phi_n : \cC_1^{(n)} \to \cC_2^{(n)}$, то $\Phi$ — эквивалентность.
\end{lemma}
\begin{proof}
Пределы сохраняют эквивалентности диаграмм; посылки восстановления и совместимости доставляют нужные диаграммы и их пределы. (Сама по себе доступность не влечёт гиперполноты или восстановления по усечениям.)
\end{proof}

\subsection{Отсутствующее обогащённое сравнение}

\begin{remark}[Теорема неограниченной стабилизации не поставляется]\label{thm:bergner-lurie-stab}
Запись $\colim_{n < \omega} (\infty, n)\text{-}\Cat$ сама по себе не доказывает перестановочности обогащения с этим копределом: семейство hom-объектов может иметь неограниченные глубины и не приходить ни с какой одной конечной стадии, а конфинальный сдвиг диаграммы не отождествляет обогащение копредела с копределом сдвинутой диаграммы. Барвик--Шоммер-Прис~\cite{BarwickSchommerPries2021} устанавливают унитарность при фиксированном конечном $n$ с пространством теорий $B(\mathbb Z/2)^n$; прежняя «теорема стабилизации Бергнер--Люри» и её доказательство перестановкой обогащения с копределом отозваны. Модель неограниченной глубины и её отображения сравнения должны задаваться явно.
\end{remark}

\section{Паранепротиворечивое расширение}\label{app:paraconsistent}

AFN-T в форме Теоремы~\ref{thm:afnt-alpha} относится к классическим Rich-метатеориям. Это приложение формулирует условное расширение на паранепротиворечивые Rich-метатеории с явно заданным классическим ядром и сохраняющей граничные предикаты трансляцией (сильно-диалетеические теории сюда не входят).

\subsection{Паранепротиворечивые R-S}

\begin{definition}[Паранепротиворечивая R-S с извлекаемым классическим ядром]\label{def:paraconsistent-rs}
Теория $S'$ — \textbf{паранепротиворечивая Rich-метатеория с извлекаемым классическим ядром}, если она удовлетворяет условиям (R1')--(R5'), аналогичным (R1)--(R5), вместе с:
\begin{itemize}
\item \textbf{(Explosion-min)}: существует формула $\phi$, такая что $\bot \not\vdash_{S'} \phi$ (эквивалентно, $\neg\forall\phi\,(\bot \vdash_{S'} \phi)$: противоречия не влекут все формулы подряд).
\item \textbf{(Strong-neg)}: существует отрицание $\sim$ такое, что $(p \wedge \sim p) \vdash_\sim \bot$ (конкретный принцип противоречия без тотального взрыва).
\item \textbf{(Извлекаемость классического ядра)}: явно заданный классический фрагмент $S^*$ непротиворечив и принадлежит $\RS$. Сохранение определений кандидатов, интерпретаций и граничных предикатов при трансляции из $S'$ составляет дополнительные данные, требуемые Теоремой~\ref{thm:paraconsistent-afnt}; автоматического существования наибольшего такого фрагмента не утверждается.
\end{itemize}
\end{definition}

\subsection{Паранепротиворечивая AFN-T}

\begin{theorem}[Условный перенос на паранепротиворечивое представление]\label{thm:paraconsistent-afnt}
Пусть заданы классическая Rich-метатеория $S^* \in \RS$ и трансляция из $S'$ в $S^*$, переводящая каждого кандидата, удовлетворяющего всем трём условиям $\LAbs$ в $S'$, в кандидата, удовлетворяющего всем трём условиям в $S^*$. Тогда в $S'$ такого кандидата не существует.
\end{theorem}
\begin{proof}[Доказательство (схема)]
Кандидат в $S'$ переводился бы трансляцией в кандидата в $S^*$, что противоречит Теореме~\ref{thm:afnt-alpha} для $S^*$.
\end{proof}
Ни сильное отрицание, ни само название «классическое ядро» не строят этой трансляции автоматически: непротиворечивость и $\RS$-принадлежность ядра, функториальность трансляции и особенно сохранение $\Pi_4$ и $\Pi_{3\text{-max}}$ остаются отдельными обязательствами, а механическое вычёркивание символа $\sim$ не является сохраняющей семантику трансляцией.

\section{Верификационный леджер Verum}\label{app:verification-ledger}

Сопутствующий \textbf{Verum MSFS Corpus} размещён по адресу \url{https://github.com/gst-st/msfs}. Его покрытие доказательствами должно сверяться с точными формулировками и допущениями настоящей ревизии: настоящий препринт не утверждает, что каждая теорема, определение или открытая конструкция имеет машинно-проверенное доказательство, и явно выделенные обязательства поднятия и стабилизации не считаются закрытыми лишь по совпадению имени метки с прежней версией.

Отчёт о верификации должен указывать ревизии корпуса и ядра, переведённую формулировку каждого пункта, все допущенные аксиомы и рамочные посылки, а также команду воспроизведения. Маркер \texttt{@framework(name, "citation")} фиксирует внешнюю зависимость, а не доказывает цитируемую теорему или её применимость. Дефинициальная консервативность, сила непротиворечивости и содержательное соответствие между формальным кодированием и текстовой формулировкой остаются самостоятельными обязательствами.

\section*{Благодарности и раскрытия}

Автор благодарит математическое сообщество за фундаментальные результаты Кантора, Рассела, Гёделя, Тарского, Фефермана, Пудлака, Ловера, Ламбека--Скотта, Адамека--Росицкого, Маккаи--Паре, Джейкобса, Пронк, Люри, Риль--Верити, Воеводского, Шрайбера и Клаузена--Шольце, на которых покоится этот синтез. Макс Середа — единственный автор; концептуализация, формальный анализ, доказательства и текст — целиком ответственность автора. Работа выполнена независимо; автор заявляет об отсутствии конкурирующих интересов и внешнего финансирования. Данные не генерировались и не анализировались; полный LaTeX-исходник заархивирован вместе с депозитом препринта.

\begin{thebibliography}{99}

\bibitem{Aczel1988}
P.~Aczel,
\newblock \emph{Non-Well-Founded Sets},
\newblock CSLI Lecture Notes 14, Stanford, 1988.

\bibitem{AdamekRosicky1994}
J.~Ad\'amek and J.~Rosick\'y,
\newblock \emph{Locally Presentable and Accessible Categories},
\newblock London Mathematical Society Lecture Notes 189, Cambridge University Press, 1994.

\bibitem{AraMaltsiniotis2017}
D.~Ara and G.~Maltsiniotis,
\newblock \emph{Homotopy comparison of cubical and globular models of weak $\omega$-categories},
\newblock Preprint, 2017.

\bibitem{AwodeyWarren2009}
S.~Awodey and M.~Warren,
\newblock Homotopy theoretic models of identity types,
\newblock \emph{Mathematical Proceedings of the Cambridge Philosophical Society} 146 (2009), 45--55.

\bibitem{Bakovic2010}
I.~Bakovi\'c,
\newblock \emph{Fibrations of bicategories},
\newblock Preprint, 2010.

\bibitem{Rezk2010}
C.~Rezk,
\newblock A Cartesian presentation of weak $n$-categories,
\newblock \emph{Geometry \& Topology} 14 (2010), no.~1, 521--571; arXiv:0901.3602.

\bibitem{Benabou1967}
J.~B\'enabou,
\newblock Introduction to bicategories,
\newblock in \emph{Reports of the Midwest Category Seminar}, Lecture Notes in Mathematics 47, Springer, 1967, pp.~1--77.

\bibitem{BergnerRezk2013}
J.E.~Bergner and C.~Rezk,
\newblock Comparison of models for $(\infty, n)$-categories, I,
\newblock \emph{Geometry \& Topology} 17 (2013), no.~4, 2163--2202.

\bibitem{BarwickSchommerPries2021}
C.~Barwick and C.~Schommer-Pries,
\newblock On the unicity of the theory of higher categories,
\newblock \emph{Journal of the American Mathematical Society} 34 (2021), 1011--1058; arXiv:1112.0040.

\bibitem{Cantor1891}
G.~Cantor,
\newblock \"Uber eine elementare Frage der Mannigfaltigkeitslehre,
\newblock \emph{Jahresbericht der Deutschen Mathematiker-Vereinigung} 1 (1891), 75--78.

\bibitem{ClausenScholze2019}
D.~Clausen and P.~Scholze,
\newblock \emph{Lectures on Condensed Mathematics},
\newblock Lecture notes, University of Bonn, 2019,
\newblock \url{https://www.math.uni-bonn.de/people/scholze/Condensed.pdf}.

\bibitem{Ernst2015}
M.~Ernst,
\newblock The prospects of unlimited category theory: doing what remains to be done,
\newblock \emph{Review of Symbolic Logic} 8 (2015), no.~2, 306--327.

\bibitem{Feferman1960}
S.~Feferman,
\newblock Arithmetization of metamathematics in a general setting,
\newblock \emph{Fundamenta Mathematicae} 49 (1960), 35--92.

\bibitem{Feferman1964}
S.~Feferman,
\newblock Systems of predicative analysis,
\newblock \emph{Journal of Symbolic Logic} 29 (1964), no.~1, 1--30.

\bibitem{Shulman2019}
M.A.~Shulman,
\newblock All $(\infty, 1)$-toposes have strict univalent universes,
\newblock preprint, 2019; arXiv:1904.07004.

\bibitem{GagnaHarpazLanari2021}
A.~Gagna, Y.~Harpaz, and E.~Lanari,
\newblock Gray tensor products and lax functors of $(\infty, 2)$-categories,
\newblock \emph{Advances in Mathematics} 391 (2021), 107986; arXiv:2002.06055.

\bibitem{GambinoGarner2008}
N.~Gambino and R.~Garner,
\newblock The identity type weak factorisation system,
\newblock \emph{Theoretical Computer Science} 409 (2008), 94--109.

\bibitem{Godel1931}
K.~G\"odel,
\newblock \"Uber formal unentscheidbare S\"atze der Principia Mathematica und verwandter Systeme I,
\newblock \emph{Monatshefte f\"ur Mathematik und Physik} 38 (1931), 173--198.

\bibitem{Hermida1999}
C.~Hermida,
\newblock Some properties of Fib as a fibred $2$-category,
\newblock \emph{Journal of Pure and Applied Algebra} 134 (1999), 83--109.

\bibitem{KapulkinLi2025}
C.~Kapulkin and Y.~Li,
\newblock Extensional concepts in intensional type theory, revisited,
\newblock \emph{Theoretical Computer Science} 1029 (2025), 115051; arXiv:2310.05706.

\bibitem{Hofmann1995}
M.~Hofmann,
\newblock \emph{Extensional concepts in intensional type theory},
\newblock PhD thesis, University of Edinburgh, 1995.

\bibitem{Hyland1982}
J.M.E.~Hyland,
\newblock The effective topos,
\newblock in \emph{The L.E.J. Brouwer Centenary Symposium}, Studies in Logic and the Foundations of Mathematics 110, North-Holland, 1982, pp.~165--216.

\bibitem{Jacobs1999}
B.~Jacobs,
\newblock \emph{Categorical Logic and Type Theory},
\newblock Studies in Logic 141, Elsevier, 1999.

\bibitem{Bishop1967}
E.~Bishop,
\newblock \emph{Foundations of Constructive Analysis},
\newblock McGraw-Hill, 1967.

\bibitem{BridgesVita2006}
D.S.~Bridges and L.S.~Vita,
\newblock \emph{Techniques of Constructive Analysis},
\newblock Universitext, Springer, 2006.

\bibitem{Buss1986}
S.R.~Buss,
\newblock \emph{Bounded Arithmetic},
\newblock Bibliopolis, Naples, 1986.

\bibitem{EseninVolpin1961}
A.S.~Esenin-Volpin,
\newblock Le programme ultra-intuitionniste des fondements des math\'ematiques,
\newblock in: \emph{Infinitistic Methods}, Pergamon Press / PWN, Warsaw, 1961, pp.~201--223.

\bibitem{EseninVolpin1970}
A.S.~Esenin-Volpin,
\newblock The ultra-intuitionistic criticism and the antitraditional program for foundations of mathematics,
\newblock in: A.~Kino, J.~Myhill, R.E.~Vesley (eds.), \emph{Intuitionism and Proof Theory}, North-Holland, 1970, pp.~3--45.

\bibitem{MarkovNagornyj1988}
A.A.~Markov and N.M.~Nagornyj,
\newblock \emph{The Theory of Algorithms},
\newblock Mathematics and its Applications (Soviet Series), Kluwer, 1988 (English translation of Nauka, Moscow, 1984).

\bibitem{Simpson2009}
S.G.~Simpson,
\newblock \emph{Subsystems of Second Order Arithmetic},
\newblock 2nd ed., Cambridge University Press, 2009.

\bibitem{LambekScott1986}
J.~Lambek and P.J.~Scott,
\newblock \emph{Introduction to Higher-Order Categorical Logic},
\newblock Cambridge Studies in Advanced Mathematics 7, Cambridge University Press, 1986.

\bibitem{KapulkinLumsdaine2021}
K.~Kapulkin and P.L.~Lumsdaine,
\newblock The simplicial model of univalent foundations (after Voevodsky),
\newblock \emph{Journal of the European Mathematical Society} 23 (2021), no.~6, 2071--2126; arXiv:1211.2851.

\bibitem{Lair1981}
C.~Lair,
\newblock Cat\'egories modelables et cat\'egories esquissables,
\newblock \emph{Diagrammes} 6 (1981), L1--L20.

\bibitem{MakkaiPare1989}
M.~Makkai and R.~Par\'e,
\newblock \emph{Accessible Categories: The Foundations of Categorical Model Theory},
\newblock Contemporary Mathematics 104, American Mathematical Society, 1989.

\bibitem{Kelly1982}
G.M.~Kelly,
\newblock \emph{Basic Concepts of Enriched Category Theory},
\newblock London Mathematical Society Lecture Notes 64, Cambridge University Press, 1982.

\bibitem{KellyStreet1974}
G.M.~Kelly and R.~Street,
\newblock Review of the elements of $2$-categories,
\newblock Lecture Notes in Mathematics 420, Springer, 1974, pp.~75--103.

\bibitem{Lawvere1969}
F.W.~Lawvere,
\newblock Adjointness in foundations,
\newblock \emph{Dialectica} 23 (1969), 281--296.

\bibitem{LurieHTT}
J.~Lurie,
\newblock \emph{Higher Topos Theory},
\newblock Annals of Mathematics Studies 170, Princeton University Press, 2009.

\bibitem{LurieHA}
J.~Lurie,
\newblock \emph{Higher Algebra},
\newblock Preprint, 2017.

\bibitem{Girard2001}
J.-Y.~Girard,
\newblock Locus solum: from the rules of logic to the logic of rules,
\newblock \emph{Mathematical Structures in Computer Science} 11 (2001), no.~3, 301--506.

\bibitem{Priest2006}
G.~Priest,
\newblock \emph{In Contradiction: A Study of the Transconsistent},
\newblock 2nd edition, Oxford University Press, 2006.

\bibitem{Pronk1996}
D.~Pronk,
\newblock Etendues and stacks as bicategories of fractions,
\newblock \emph{Compositio Mathematica} 102 (1996), no.~3, 243--303.

\bibitem{Pudlak1985}
P.~Pudl\'ak,
\newblock Cuts, consistency statements and interpretations,
\newblock \emph{Journal of Symbolic Logic} 50 (1985), no.~2, 423--441.

\bibitem{Rathjen1999}
M.~Rathjen,
\newblock The realm of ordinal analysis,
\newblock in \emph{Sets and Proofs} (Leeds 1997), London Mathematical Society Lecture Notes 258, Cambridge University Press, 1999, pp.~219--279.

\bibitem{RiehlVerity2022}
E.~Riehl and D.~Verity,
\newblock \emph{Elements of $\infty$-Category Theory},
\newblock Cambridge Studies in Advanced Mathematics 194, Cambridge University Press, 2022.

\bibitem{Russell1903}
B.~Russell,
\newblock \emph{The Principles of Mathematics},
\newblock Cambridge University Press, 1903.

\bibitem{Schreiber2013}
U.~Schreiber,
\newblock \emph{Differential cohomology in a cohesive $\infty$-topos},
\newblock arXiv:1310.7930, 2013.

\bibitem{Shulman2008}
M.A.~Shulman,
\newblock $2$-toposes,
\newblock \emph{Journal of Pure and Applied Algebra} 210 (2008), 273--287.

\bibitem{Streicher1991}
T.~Streicher,
\newblock \emph{Semantics of type theory},
\newblock Birkh\"auser Progress in Theoretical Computer Science, 1991.

\bibitem{Tarski1936}
A.~Tarski,
\newblock Der Wahrheitsbegriff in den formalisierten Sprachen,
\newblock \emph{Studia Philosophica} 1 (1936), 261--405.

\bibitem{HoTTBook2013}
The Univalent Foundations Program,
\newblock \emph{Homotopy Type Theory: Univalent Foundations of Mathematics},
\newblock Institute for Advanced Study, Princeton, 2013,
\newblock \url{https://homotopytypetheory.org/book/}.

\bibitem{VoevodskyFoundations2010}
V.~Voevodsky,
\newblock \emph{Foundations} Coq library,
\newblock software artefact, 2010--2015,
\newblock \url{https://github.com/vladimirias/Foundations}.

\bibitem{Werner1994}
B.~Werner,
\newblock \emph{Une Th\'eorie des Constructions Inductives},
\newblock PhD thesis, Universit\'e Paris 7, 1994.

\bibitem{Yanofsky2003}
N.S.~Yanofsky,
\newblock A universal approach to self-referential paradoxes, incompleteness and fixed points,
\newblock \emph{Bulletin of Symbolic Logic} 9 (2003), 362--386.

\end{thebibliography}

\end{document}
